% Generated mechanically from the frozen stacks-geometry.tex target.
% Do not edit this generated file; edit the bound locale source instead.

% OK, start here.
%

\setcounter{chapter}{106}
\chapter{代数叠的几何}



\phantomsection
\label{stacks-geometry-section-phantom}


\section{引言}
\label{stacks-geometry-section-introduction}

\noindent
本章讨论代数叠的若干几何性质。第
\ref{stacks-geometry-section-multiplicities} 节和第
\ref{stacks-geometry-section-dimension-of-algebraic-stacks} 节的初始版本由
Matthew Emerton 与 Toby Gee 撰写，其原始形式见
\cite{Emerton-Gee-dim}。







\section{半通用环}
\label{stacks-geometry-section-versal}

\noindent
本节阐明形变环与局部 Noether 基上有限型代数叠的局部环之间的关系。

\begin{situation}
\label{stacks-geometry-situation-versal}
这里 $\mathcal{X}$ 是局部 Noether 概形 $S$ 上的局部有限型代数叠。
\end{situation}

\noindent
定义如下。

\begin{definition}
\label{stacks-geometry-definition-versal-ring-at-x}
在情形 \ref{stacks-geometry-situation-versal} 中，设
$x_0 : \Spec(k) \to \mathcal{X}$ 为一态射，其中 $k$ 是 $S$ 上的
有限型域。$\mathcal{X}$ 在 $x_0$ 处的一个
{\it 半通用环（versal ring）}，是一个剩余域为 $k$ 的完备
Noether 局部 $S$-代数 $A$，使得存在 Artin 公理中定义
\ref{artin-definition-versal-formal-object} 所述的半通用形式对象
$(A, \xi_n, f_n)$，且 $\xi_1 \cong x_0$（一个 $2$-同构）。
\end{definition}

\noindent
我们要证明半通用环存在，并且在光滑因子意义下唯一。为此将使用
《Artin 公理》第 \ref{artin-section-predeformation-categories} 节中的
预形变范畴。在我们的情形中，它们总是形变范畴。

\begin{lemma}
\label{stacks-geometry-lemma-deformation-category}
在情形 \ref{stacks-geometry-situation-versal} 中，设
$x_0 : \Spec(k) \to \mathcal{X}$ 为一态射，其中 $k$ 是 $S$ 上的
有限型域。那么 $\mathcal{F}_{\mathcal{X}, k, x_0}$
是形变范畴，并且
$T\mathcal{F}_{\mathcal{X}, k, x_0}$ 与
$\text{Inf}(\mathcal{F}_{\mathcal{X}, k, x_0})$
都是有限维 $k$-向量空间。
\end{lemma}

\begin{proof}
取仿射开子概形 $\Spec(\Lambda) \subset S$，使得
$\Spec(k) \to S$ 经它分解。由《Artin 公理》第
\ref{artin-section-predeformation-categories} 节，得到范畴
$\mathcal{C}_\Lambda$ 上的预形变范畴
$\mathcal{F}_{\mathcal{X}, k, x_0}$。
（正如所引之处指出的，该范畴只依赖于态射
$\Spec(k) \to S$，而不依赖于 $\Lambda$ 的选择。）
由《Artin 公理》中的引理
\ref{artin-lemma-deformation-category} 和
\ref{artin-lemma-algebraic-stack-RS}，
$\mathcal{F}_{\mathcal{X}, k, x_0}$ 实际上是形变范畴。
再由《Artin 公理》中的引理 \ref{artin-lemma-finite-dimension}，
$T\mathcal{F}_{\mathcal{X}, k, x_0}$ 与
$\text{Inf}(\mathcal{F}_{\mathcal{X}, k, x_0})$
是有限维 $k$-向量空间。
\end{proof}




\begin{lemma}
\label{stacks-geometry-lemma-versal-ring}
在情形 \ref{stacks-geometry-situation-versal} 中，设
$x_0 : \Spec(k) \to \mathcal{X}$ 为一态射，其中 $k$ 是 $S$ 上的
有限型域。那么 $\mathcal{X}$ 在 $x_0$ 处存在半通用环。
给定其中任意一对 $A$、$A'$，则对某个 $r$，作为 $S$-代数有
$A \cong A'[[t_1, \ldots, t_r]]$ 或
$A' \cong A[[t_1, \ldots, t_r]]$。
\end{lemma}

\begin{proof}
存在性由引理 \ref{stacks-geometry-lemma-deformation-category} 以及《形式形变论》中的
定义 \ref{formal-defos-definition-deformation-category} 和引理
\ref{formal-defos-lemma-RS-implies-S1-S2}、
\ref{formal-defos-lemma-versal-object-existence} 得出。
由《形式形变论》中的唯一性结果——引理
\ref{formal-defos-lemma-minimal-versal}，存在 $\mathcal{X}$ 在
$x_0$ 处的一个“极小”半通用环 $A$，使得 $\mathcal{X}$ 在 $x_0$
处的任何其他半通用环都同构于某个
$A[[t_1, \ldots, t_r]]$，其中 $r$ 适当。这显然推出第二个陈述。
\end{proof}

\begin{lemma}
\label{stacks-geometry-lemma-versal-ring-field-extension}
在情形 \ref{stacks-geometry-situation-versal} 中，设
$x_0 : \Spec(k) \to \mathcal{X}$ 为一态射，其中 $k$ 是 $S$ 上的
有限型域。设 $l/k$ 是有限域扩张，并用
% P12-ERR-0145: see ../control/ERRATA_CANDIDATES.jsonl
$x_{l, 0} : \Spec(l) \to \mathcal{X}$ 表示诱导的态射。
给定 $\mathcal{X}$ 在 $x_0$ 处的一个半通用环 $A$，存在
$\mathcal{X}$ 在 $x_{l, 0}$ 处的半通用环 $A'$，使得有一个
$S$-代数映射 $A \to A'$，它诱导给定的域扩张 $l/k$，并且在
$\mathfrak m_{A'}$-adic 拓扑下形式光滑。
\end{lemma}

\begin{proof}
这立即由《Artin 公理》中的引理
\ref{artin-lemma-change-of-field} 和《形式形变论》中的引理
\ref{formal-defos-lemma-change-of-field-versal-ring} 得出。
（还使用了 $\mathcal{X}$ 由《Artin 公理》中的引理
\ref{artin-lemma-algebraic-stack-RS} 满足 (RS)。）
\end{proof}

\begin{lemma}
\label{stacks-geometry-lemma-compare-versal-ring-completion}
在情形 \ref{stacks-geometry-situation-versal} 中，设 $x : U \to \mathcal{X}$ 为一态射，
其中 $U$ 是 $S$ 上的局部有限型概形。设 $u_0 \in U$ 是有限型点。
令 $k = \kappa(u_0)$，并用
% P12-ERR-0146: see ../control/ERRATA_CANDIDATES.jsonl
$x_0 : \Spec(k) \to \mathcal{X}$ 表示诱导的映射。下列条件等价：
\begin{enumerate}
\item $x$ 在 $u_0$ 处半通用
（《Artin 公理》中的定义 \ref{artin-definition-versal}）；
\item
$\hat x : \mathcal{F}_{U, k, u_0} \to \mathcal{F}_{\mathcal{X}, k, x_0}$
光滑；
\item 与
$x|_{\Spec(\mathcal{O}_{U, u_0}^\wedge)}$ 相关联的形式对象半通用；以及
% P12-ERR-0147: see ../control/ERRATA_CANDIDATES.jsonl
\item 存在 $U' \subset U$，它是 $x$ 的开邻域，并且
$x|_{U'} : U' \to \mathcal{X}$ 光滑。
\end{enumerate}
而且在此情形下，完备化 $\mathcal{O}_{U, u_0}^\wedge$ 是
$\mathcal{X}$ 在 $x_0$ 处的半通用环。
\end{lemma}

\begin{proof}
由于 $U \to S$ 局部有限型（它是这类态射的复合），
$\Spec(k) \to S$ 也是有限型的（同样是因为它是这类态射的复合）。
因此该陈述有意义。(1) 与 (2) 的等价性正是 $x$ 在 $u_0$ 处半通用
的定义。(1) 与 (3) 的等价性由《Artin 公理》中的引理
\ref{artin-lemma-versality-matches} 给出。因此 (1)、(2)、(3) 等价。

\medskip\noindent
若 $x|_{U'}$ 光滑，则由《Artin 公理》中的引理
\ref{artin-lemma-formally-smooth-on-deformation-categories}，函子
$\hat x : \mathcal{F}_{U, k, u_0} \to \mathcal{F}_{\mathcal{X}, k, x_0}$
光滑。因此 (4) 推出 (1)、(2)、(3)。
反过来，设 $x$ 在 $u_0$ 处半通用。取满的光滑态射
$y : V \to \mathcal{X}$，其中 $V$ 是概形。令
$Z = V \times_\mathcal{X} U$，并取位于 $u_0$ 上方的有限型点
$z_0 \in |Z|$（这是可能的，参见《空间态射》中的引理
\ref{spaces-morphisms-lemma-finite-type-points-surjective-morphism}）。
由《Artin 公理》中的引理 \ref{artin-lemma-base-change-versal}，
态射 $Z \to V$ 在 $z_0$ 处光滑。按定义可找到 $z_0$ 的开邻域
$W \subset Z$，使得 $W \to V$ 光滑。由于 $Z \to U$ 是开映射，
令 $U' \subset U$ 为 $W$ 的像。于是根据我们对叠的光滑态射的定义，
$U' \to \mathcal{X}$ 光滑。

\medskip\noindent
最后一个陈述由定义得出，因为 $\mathcal{O}_{U, u_0}^\wedge$
pro-表示 $\mathcal{F}_{U, k, u_0}$。
\end{proof}

\begin{lemma}
\label{stacks-geometry-lemma-characterize-smoothness}
% P12-ERR-0148: see ../control/ERRATA_CANDIDATES.jsonl
在情形 \ref{stacks-geometry-situation-versal} 中，设
$x_0 : \Spec(k) \to \mathcal{X}$ 为一态射，使得
$\Spec(k) \to S$ 是有限型的，其像为 $s$。设 $A$ 是
$\mathcal{X}$ 在 $x_0$ 处的半通用环。下列条件等价：
\begin{enumerate}
\item $x_0$ 位于 $\mathcal{X} \to S$ 的光滑轨迹中
（《叠的态射》中的引理 \ref{stacks-morphisms-lemma-where-smooth}）；
\item $\mathcal{O}_{S, s} \to A$ 在
$\mathfrak m_A$-adic 拓扑下形式光滑；以及
\item $\mathcal{F}_{\mathcal{X}, k, x_0}$ 无阻碍。
\end{enumerate}
\end{lemma}

\begin{proof}
(2) 与 (3) 的等价性立即由《形式形变论》中的引理
\ref{formal-defos-lemma-smooth-power-series-classical} 得出。

\medskip\noindent
注意，$\mathcal{O}_{S, s} \to A$ 在
$\mathfrak m_A$-adic 拓扑下形式光滑，当且仅当
$\mathcal{O}_{S, s} \to A' = A[[t_1, \ldots, t_r]]$
在 $\mathfrak m_{A'}$-adic 拓扑下形式光滑。
因此由引理 \ref{stacks-geometry-lemma-versal-ring}，(2) 不依赖于半通用环的选择。
接着，设 $l/k$ 是有限扩张，并按引理
\ref{stacks-geometry-lemma-versal-ring-field-extension} 选择 $A \to A'$。
若 $\mathcal{O}_{S, s} \to A$ 在 $\mathfrak m_A$-adic 拓扑下
形式光滑，则 $\mathcal{O}_{S, s} \to A'$ 在
$\mathfrak m_{A'}$-adic 拓扑下形式光滑；参见《代数的进一步讨论》
中的引理 \ref{more-algebra-lemma-compose-formally-smooth}。
反过来，若 $\mathcal{O}_{S, s} \to A'$ 在
$\mathfrak m_{A'}$-adic 拓扑下形式光滑，则
$\mathcal{O}_{S, s}^\wedge \to A'$ 与 $A \to A'$ 正则
（《代数的进一步讨论》中的命题
\ref{more-algebra-proposition-fs-regular}），从而
$\mathcal{O}_{S, s}^\wedge \to A$ 正则
（《代数的进一步讨论》中的引理
\ref{more-algebra-lemma-regular-permanence}），所以
$\mathcal{O}_{S, s} \to A$ 在 $\mathfrak m_A$-adic 拓扑下形式光滑
（仍用前一引理）。因此，对 $k$ 和 $x_0$，(2) 与 (1) 等价，
当且仅当对 $l$ 和 $x_{0, l}$ 也等价。

\medskip\noindent
取概形 $U$ 和光滑态射 $U \to \mathcal{X}$，使得
$\Spec(k) \times_\mathcal{X} U$ 非空。取有限扩张 $l/k$ 和一点
$w_0 : \Spec(l) \to \Spec(k) \times_\mathcal{X} U$。
设 $u_0 \in U$ 为 $w_0$ 的像。将上面的论证分别应用于
$l/k$ 与 $l/\kappa(u_0)$，可知能够约化到 $u_0$。因此可以设
$A = \mathcal{O}_{U, u_0}^\wedge$；参见引理
\ref{stacks-geometry-lemma-compare-versal-ring-completion}。
注意，$x_0$ 位于 $\mathcal{X} \to S$ 的光滑轨迹中，当且仅当
$u_0$ 位于 $U \to S$ 的光滑轨迹中；例如参见《叠的态射》中的引理
\ref{stacks-morphisms-lemma-where-smooth}。
于是 (1) 与 (2) 的等价性由《代数的进一步讨论》中的引理
\ref{more-algebra-lemma-formally-smooth-finite-type} 得出。
\end{proof}

\noindent
回顾 Artin 逼近的一个推论。

\begin{lemma}
\label{stacks-geometry-lemma-Artin-approximation-by-smooth-morphism}
% P12-ERR-0148: see ../control/ERRATA_CANDIDATES.jsonl
在情形 \ref{stacks-geometry-situation-versal} 中，设
$x_0 : \Spec(k) \to \mathcal{X}$ 为一态射，使得
$\Spec(k) \to S$ 是有限型的，其像为 $s$。设 $A$ 是
$\mathcal{X}$ 在 $x_0$ 处的半通用环。
若 $\mathcal{O}_{S, s}$ 是 G-环，则可找到一个光滑态射
$U \to \mathcal{X}$，其源是概形，并可找到一个剩余域为 $k$ 的点
$u_0 \in U$，使得
\begin{enumerate}
\item $\Spec(k) \to U \to \mathcal{X}$ 与给定态射 $x_0$ 相同；
\item 存在同构 $\mathcal{O}_{U, u_0}^\wedge \cong A$。
\end{enumerate}
\end{lemma}

\begin{proof}
设 $(\xi_n, f_n)$ 是 $A$ 上的半通用形式对象。
由《Artin 公理》中的引理 \ref{artin-lemma-effective}，
$\xi = (A, \xi_n, f_n)$ 是有效的。
根据假设，$\mathcal{X}$ 在 $S$ 上局部有限表示
（使用《叠的态射》中的引理
\ref{stacks-morphisms-lemma-noetherian-finite-type-finite-presentation}），
因而由《叠的极限》中的命题
\ref{stacks-limits-proposition-characterize-locally-finite-presentation}
保极限。因此，按《Artin 公理》中的引理
\ref{artin-lemma-approximate-versal} 所述的 Artin 逼近表明，
可以找到态射 $U \to \mathcal{X}$，其源为有限型 $S$-概形，
并含有剩余域为 $k$ 的一点 $u_0 \in U$，它满足 (1) 与 (2)，
且 $U \to \mathcal{X}$ 在 $u_0$ 处半通用。
由引理 \ref{stacks-geometry-lemma-compare-versal-ring-completion}，缩小 $U$ 后，
可以设 $U \to \mathcal{X}$ 光滑。
\end{proof}

\begin{remark}[提升半通用环]
\label{stacks-geometry-remark-upgrade}
在情形 \ref{stacks-geometry-situation-versal} 中，设
$x_0 : \Spec(k) \to \mathcal{X}$ 为一态射，其中 $k$ 是 $S$ 上的
有限型域。设 $A$ 是 $\mathcal{X}$ 在 $x_0$ 处的半通用环。
由《Artin 公理》中的引理 \ref{artin-lemma-effective}，
我们的半通用形式对象实际上来自 $S$ 上的一个态射
$$
\Spec(A) \longrightarrow \mathcal{X}
$$
此外，上述每个结果都可以提升为与该态射相容。列举如下：
% P12-ERR-0149: see ../control/ERRATA_CANDIDATES.jsonl
\begin{enumerate}
\item 在引理 \ref{stacks-geometry-lemma-versal-ring} 中，同构
$A \cong A'[[t_1, \ldots, t_r]]$ 或
$A' \cong A[[t_1, \ldots, t_r]]$
可以选成与这些态射相容；
\item 在引理 \ref{stacks-geometry-lemma-versal-ring-field-extension} 中，
同态 $A \to A'$ 可以选成与这些态射相容；
\item 在引理 \ref{stacks-geometry-lemma-compare-versal-ring-completion} 中，
态射 $\Spec(\mathcal{O}_{U, u_0}^\wedge) \to \mathcal{X}$
是典范映射 $\Spec(\mathcal{O}_{U, u_0}^\wedge) \to U$
与给定映射 $U \to \mathcal{X}$ 的复合；
\item 在引理 \ref{stacks-geometry-lemma-Artin-approximation-by-smooth-morphism} 中，
同构 $\mathcal{O}_{U, u_0}^\wedge \cong A$ 可以选成使
$\Spec(A) \to \mathcal{X}$ 对应于上一项中的典范映射。
\end{enumerate}
在每种情形中，该陈述都由我们的映射与半通用形式元素相容这一事实
得出；但要注意，所蕴含的图表仅在对一个 $2$-箭头作
（非典范）选择后才是 $2$-交换的。尽管如此，这意味着 (1) 中
所蕴含的映射 $A' \to A$ 或 $A \to A'$ 在形式同伦意义下是
良定义的；参见《形式形变论》中的引理
\ref{formal-defos-lemma-versal-unique-up-to-homotopy}。
\end{remark}

\begin{lemma}
\label{stacks-geometry-lemma-versal-ring-flat}
在情形 \ref{stacks-geometry-situation-versal} 中，设
$x_0 : \Spec(k) \to \mathcal{X}$ 为一态射，其中 $k$ 是 $S$ 上的
有限型域。设 $A$ 是 $\mathcal{X}$ 在 $x_0$ 处的半通用环。
那么注 \ref{stacks-geometry-remark-upgrade} 中的态射
$\Spec(A) \to \mathcal{X}$ 是平坦的。
\end{lemma}

\begin{proof}
如果 $S$ 在像点处的局部环是 G-环，则这立即由引理
\ref{stacks-geometry-lemma-Artin-approximation-by-smooth-morphism} 以及
Noether 局部环到其完备化的映射平坦这一事实得出。一般情形证明如下。

\medskip\noindent
步骤 I。若 $A$ 和 $A'$ 是 $\mathcal{X}$ 在 $x_0$ 处的两个
半通用环，则该结果对 $A$ 成立，当且仅当对 $A'$ 成立。
事实上，必要时交换 $A$ 与 $A'$ 后，可以设存在形式光滑映射
$\varphi : A \to A'$，使得复合
$$
\Spec(A') \to \Spec(A) \to \mathcal{X}
$$
就是态射 $\Spec(A') \to \mathcal{X}$；参见引理
\ref{stacks-geometry-lemma-versal-ring} 与注 \ref{stacks-geometry-remark-upgrade}。
由于 $A \to A'$ 忠实平坦，所需等价性由《叠的态射》中的引理
\ref{stacks-morphisms-lemma-composition-flat} 和
\ref{stacks-morphisms-lemma-flat-permanence} 得出。

\medskip\noindent
步骤 II。设 $l/k$ 是有限域扩张。设
$x_{l, 0} : \Spec(l) \to \mathcal{X}$ 为诱导的态射。
设 $A$ 是 $\mathcal{X}$ 在 $x_0$ 处的半通用环，并设
$A \to A'$ 如引理 \ref{stacks-geometry-lemma-versal-ring-field-extension} 中所述。
此时复合
$$
\Spec(A') \to \Spec(A) \to \mathcal{X}
$$
仍然是态射 $\Spec(A') \to \mathcal{X}$；参见注
\ref{stacks-geometry-remark-upgrade}。与之前相同地论证，并用步骤 I 说明半通用环的
选择无关紧要，就得到该引理对 $x_0$ 成立，当且仅当对
$x_{l, 0}$ 成立。

\medskip\noindent
步骤 III。取概形 $U$ 和满的光滑态射
$U \to \mathcal{X}$。于是可以在
$Z = U \times_\mathcal{X} x_0$（这是非空代数空间）上取有限型点
$z_0$。设 $u_0 \in U$ 为 $z_0$ 在 $U$ 中的像。
取一个概形和满的 \'etale 映射 $W \to Z$，使得 $z_0$ 是
闭点 $w_0 \in W$ 的像（参见《空间态射》第
\ref{spaces-morphisms-section-points-finite-type} 节）。
由于 $W \to \Spec(k)$ 和 $W \to U$ 都是有限型的，
$\kappa(w_0)/k$ 与 $\kappa(w_0)/\kappa(u_0)$ 都是有限域扩张
（参见《态射》第 \ref{morphisms-section-points-finite-type} 节）。
两次应用步骤 II，可以把 $x_0$ 替换为
$u_0 \to U \to \mathcal{X}$。于是我们的态射是复合
$$
\Spec(\mathcal{O}_{U, u_0}^\wedge) \to U \to \mathcal{X}
$$
% P12-ERR-0150: see ../control/ERRATA_CANDIDATES.jsonl
第一支箭头平坦，因为 Noether 局部环的完备化平坦
（《代数》中的引理 \ref{algebra-lemma-completion-flat}）；
第二支箭头平坦，因为光滑态射平坦。复合平坦，因为复合保持平坦性。
\end{proof}

\begin{remark}
\label{stacks-geometry-remark-groupoid-defo}
在情形 \ref{stacks-geometry-situation-versal} 中，设
$x_0 : \Spec(k) \to \mathcal{X}$ 为一态射，其中 $k$ 是 $S$ 上的
有限型域。由引理 \ref{stacks-geometry-lemma-deformation-category} 和《形式形变论》
中的定理 \ref{formal-defos-theorem-presentation-deformation-groupoid}，
可知 $\mathcal{F}_{\mathcal{X}, k, x_0}$ 在
$\mathcal{C}_\Lambda$ 上的函子中有一个由光滑 pro-可表群胚给出的
表示。展开定义，这意味着可以选择
\begin{enumerate}
\item 剩余域为 $k$ 的 Noether 完备局部 $\Lambda$-代数 $A$，
以及 $A$ 上的 $\mathcal{F}_{\mathcal{X}, k, x_0}$ 的半通用形式对象
$\xi$；
\item 剩余域为 $k$ 的 Noether 完备局部 $\Lambda$-代数 $B$，
以及同构
$$
\underline{B}|_{\mathcal{C}_\Lambda}
\longrightarrow
\underline{A}|_{\mathcal{C}_\Lambda}
\times_{\underline{\xi}, \mathcal{F}_{\mathcal{X}, k, x_0}, \underline{\xi}}
\underline{A}|_{\mathcal{C}_\Lambda}
$$
\end{enumerate}
两个投影对应于形式光滑映射
$t : A \to B$ 和 $s : A \to B$（因为 $\xi$ 半通用）。
存在映射 $c : B \to B \widehat{\otimes}_{s, A, t} B$，
它使 $(A, B, s, t, c)$ 成为剩余域为 $k$ 的 Noether 完备局部
$\Lambda$-代数范畴中的余群胚
（在 pro-可表函子上，该映射由《形式形变论》中的引理
\ref{formal-defos-lemma-presentation-construction} 构造）。
最后，所引定理说明 $\xi$ 诱导等价
$$
[\underline{A}|_{\mathcal{C}_\Lambda} / \underline{B}|_{\mathcal{C}_\Lambda}]
\longrightarrow
\mathcal{F}_{\mathcal{X}, k, x_0}
$$
它是 $\mathcal{C}_\Lambda$ 上余纤维化群胚之间的等价。
事实上，还得到等价
$$
[\underline{A}/\underline{B}]
\longrightarrow
\widehat{\mathcal{F}}_{\mathcal{X}, k, x_0}
$$
它是完备范畴 $\widehat{\mathcal{C}}_\Lambda$ 上余纤维化群胚之间的
等价（关于这为何成立，参见《形式形变论》第
\ref{formal-defos-section-prorepresentable-groupoids-in-functors} 节
中的讨论）。当然，$A$ 是 $\mathcal{X}$ 在 $x_0$ 处的半通用环。
\end{remark}








\section{代数叠各不可约分支的重数}
\label{stacks-geometry-section-multiplicities}

\noindent
若 $X$ 是局部 Noether 概形，则可将 $X$（仅视为拓扑空间）写成
不可约分支的并，记为 $X = \bigcup T_i$。每个不可约分支都是唯一
一般点 $\xi_i$ 的闭包，而局部环 $\mathcal O_{X,\xi_i}$ 是局部
Artin 环。可以把 {\it $X$ 沿 $T_i$ 的重数}，或称
{\it $T_i$ 在 $X$ 中的重数}，定义为
$$
m_{T_i, X} = \text{长度}_{\mathcal O_{X, \xi_i}} \mathcal O_{X, \xi_i}
$$
换言之，它就是该局部 Artin 环的长度。请与《Chow 同调》第
\ref{chow-section-cycle-of-closed-subscheme} 节比较。

\medskip\noindent
这里的目标是把该定义推广到局部 Noether 代数叠。若 $\mathcal{X}$
是叠，则其拓扑空间 $|\mathcal{X}|$（见《叠的性质》中的定义
\ref{stacks-properties-definition-topological-space}）是局部 Noether 的
（《叠的态射》中的引理
\ref{stacks-morphisms-lemma-Noetherian-topology}）。
$|\mathcal{X}|$ 的不可约分支有时也称为 $\mathcal{X}$ 的不可约分支。
若 $\mathcal{X}$ 拟分离，则 $|\mathcal{X}|$ 是清醒的（sober）
（《叠的态射》中的引理 \ref{stacks-morphisms-lemma-sober-qs}），
但在非拟分离情形下未必如此。例如可考虑非拟分离代数空间
$X = \mathbf{A}^1_\mathbf{C}/\mathbf{Z}$。此外，$|\mathcal{X}|$ 上
并不存在一个可以利用其茎来定义重数的结构层。

\begin{lemma}
\label{stacks-geometry-lemma-map-of-components}
设 $f : U \to \mathcal{X}$ 是从概形到局部 Noether 代数叠的光滑态射。
$|U|$ 的任一不可约分支之像的闭包，都是 $|\mathcal{X}|$ 的不可约分支。
若 $U \to \mathcal{X}$ 满射，则 $|\mathcal{X}|$ 的所有不可约分支都可
由此得到。
\end{lemma}

\begin{proof}
由《叠的性质》中的引理 \ref{stacks-properties-lemma-topology-points}，
映射 $|U| \to |\mathcal{X}|$ 连续且开。设 $T \subset |U|$ 是一个
不可约分支。由于 $U$ 局部 Noether，可以找到包含于 $T$ 的非空仿射开集
$W \subset U$。于是 $f(T) \subset |\mathcal{X}|$ 不可约，并且包含
非空开子集 $f(W)$。因此 $f(T)$ 的闭包不可约且包含一个非空开集，
从而该闭包是不可约分支。

\medskip\noindent
假设 $U \to \mathcal{X}$ 满射，并设 $Z \subset |\mathcal{X}|$ 是
不可约分支。取 $|\mathcal{X}|$ 的一个与 $Z$ 相交的 Noether 开子集
$V$。从 $V$ 中删去其他不可约分支后，可以设 $V \subset Z$。
取非空开集 $f^{-1}(V) \subset |U|$ 的一个不可约分支，并设
$T \subset |U|$ 为其闭包。它是 $|U|$ 的不可约分支；由 $T$ 的选择，
$f(T)$ 的闭包必与 $Z$ 相同。
\end{proof}

\noindent
上一引理尤其适用于局部 Noether 概形之间的光滑态射。下一个引理的
陈述中将隐含地使用这一特例。

\begin{lemma}
\label{stacks-geometry-lemma-multiplicities}
设 $U \to X$ 是局部 Noether 概形之间的光滑态射。
% P12-ERR-0151: see ../control/ERRATA_CANDIDATES.jsonl
设 $T'$ 是 $U$ 的一个不可约分支。设 $T$ 是 $T'$ 的像的闭包所给出的
$X$ 的不可约分支。那么 $m_{T', U} = m_{T, X}$。
\end{lemma}

\begin{proof}
用 $\xi'$ 表示 $T'$ 的一般点，用 $\xi$ 表示 $T$ 的一般点。
设 $A = \mathcal{O}_{X, \xi}$ 且 $B = \mathcal{O}_{U, \xi'}$。
需要证明 $\text{长度}_A A = \text{长度}_B B$。由于
$A \to B$ 是平坦的局部环同态（因为光滑态射平坦），由《代数》中的
引理 \ref{algebra-lemma-pullback-module} 有
$$
\text{长度}_A(A) \text{长度}_B(B/\mathfrak m_A B) =
\text{长度}_B(B)
$$
因此只需证明 $\mathfrak m_A B = \mathfrak m_B$，等价地，只需证明
$B/\mathfrak m_A B$ 约化。由于 $U \to X$ 光滑，其基变换
$U_{\xi} \to \Spec \kappa(\xi)$ 也光滑。由于 $U_{\xi}$ 是域上的
光滑概形，它是约化的，因而其任一点处的局部环也约化
% P12-ERR-0152: see ../control/ERRATA_CANDIDATES.jsonl
（《簇》中的引理 \ref{varieties-lemma-smooth-geometrically-normal}）。特别地，
$$
B/\mathfrak m_A B =
\mathcal{O}_{U, \xi'}/\mathfrak m_{X, \xi}\mathcal{O}_{U, \xi'} =
\mathcal{O}_{U_\xi, \xi'}
$$
约化，这正是所需。
\end{proof}

\noindent
利用这一结果，可以通过在光滑局部考察来说明重数有一个良定义的概念。

\begin{lemma}
\label{stacks-geometry-lemma-multiplicity}
设 $U_1 \to \mathcal{X}$ 和 $U_2 \to \mathcal{X}$ 是从概形到局部
Noether 代数叠 $\mathcal{X}$ 的两个光滑态射。分别设 $T_1'$ 和
$T_2'$ 是 $|U_1|$ 与 $|U_2|$ 的不可约分支。假设 $T_1'$ 与
$T_2'$ 的像的闭包，是 $|\mathcal{X}|$ 的同一个不可约分支 $T$。
那么 $m_{T_1', U_1} = m_{T_2', U_2}$。
\end{lemma}

\begin{proof}
分别设 $V_1$ 与 $V_2$ 是 $T_1'$ 与 $T'_2$ 的稠密子集，并分别在
$U_1$ 与 $U_2$ 中开（见引理 \ref{stacks-geometry-lemma-map-of-components} 的证明）。
$|V_1|$ 与 $|V_2|$ 在 $|\mathcal{X}|$ 中的像，都是不可约子集 $T$
的非空开子集，因而有非空交。由《叠的性质》中的引理
\ref{stacks-properties-lemma-points-cartesian}，映射
$|V_1 \times_\mathcal{X} V_2| \to |V_1| \times_{|\mathcal{X}|} |V_2|$
满射。因此 $V_1 \times_\mathcal{X} V_2$ 是非空代数空间；于是可取
一个源为（非空）概形的 \'etale 满射
$V \to V_1 \times_\mathcal{X} V_2$。
若以 $T'$ 表示 $V$ 的任一不可约分支，则引理
\ref{stacks-geometry-lemma-map-of-components} 表明，$T'$ 在 $U_1$ 中的像的闭包
（相应地，在 $U_2$ 中的像的闭包）等于 $T'_1$（相应地，等于 $T'_2$）。

\medskip\noindent
两次应用引理 \ref{stacks-geometry-lemma-multiplicities}，得到
$$
m_{T_1', U_1} = m_{T', V} = m_{T_2', U_2},
$$
这正是所需。
\end{proof}

\noindent
至此已经完成足够的工作，从而说明下述定义有意义。

\begin{definition}
\label{stacks-geometry-definition-multiplicity}
设 $\mathcal{X}$ 是局部 Noether 代数叠。设
$T \subset |\mathcal{X}|$ 是一个不可约分支。$T$ 在 $\mathcal{X}$ 中的
{\it 重数}定义为 $m_{T, \mathcal{X}} = m_{T', U}$，其中
$f : U \to \mathcal{X}$ 是从概形出发的光滑态射，而
$T' \subset |U|$ 是满足 $f(T') \subset T$ 的不可约分支。
\end{definition}

\noindent
由引理 \ref{stacks-geometry-lemma-map-of-components} 和 \ref{stacks-geometry-lemma-multiplicity}，
该定义与 $f : U \to \mathcal{X}$ 的选择以及映到 $T$ 的不可约分支
$T'$ 的选择无关。

\medskip\noindent
最后指出，把 $\mathcal{X}$ 的不可约分支视为闭子叠有时很方便。
为此，若 $T \subset |\mathcal{X}|$ 是不可约分支，则可考虑唯一的
约化闭子叠 $\mathcal{T} \subset \mathcal{X}$，使得
$|\mathcal{T}| = T$；见《叠的性质》中的定义
\ref{stacks-properties-definition-reduced-induced-stack}。
若 $\mathcal{X}$ 拟分离，则不可约分支是整叠；进一步讨论见
《叠的态射》第 \ref{stacks-morphisms-section-integral-stacks} 节。



\section{形式分支与重数}
\label{stacks-geometry-section-formal-branches}

\noindent
定义 \ref{stacks-geometry-definition-multiplicity} 给出的不可约分支重数概念，与
$\mathcal{X}$ 在有限型点处的半通用环（之谱）的不可约分支重数这一
相关概念之间，有一个比较将很方便。

\medskip\noindent
在情形 \ref{stacks-geometry-situation-versal} 中，设 $x_0 : \Spec(k) \to \mathcal{X}$
为一态射，其中 $k$ 是 $S$ 上的有限型域。设 $A$、$A'$ 是
$\mathcal{X}$ 在 $x_0$ 处的半通用环。必要时交换 $A$ 与 $A'$ 后，
可知存在一个形式光滑的\footnote{意指 $A'$ 同构于 $A$ 上的幂级数环。}
映射 $\varphi : A \to A'$，它与半通用形式对象相容；见引理
\ref{stacks-geometry-lemma-versal-ring} 和注 \ref{stacks-geometry-remark-upgrade}。
此外，$\varphi$ 在形式同伦意义下良定义；见《形式形变论》中的引理
\ref{formal-defos-lemma-versal-unique-up-to-homotopy}。特别地，由
《形式形变论》中的引理
\ref{formal-defos-lemma-homotopic-minimal-prime}，
$\varphi(\mathfrak p)A'$ 是 $A'$ 的良定义理想。由于 $A \to A'$
形式光滑，事实上 $\varphi(\mathfrak p)A'$ 是 $A'$ 的极小素理想，
而 $A'$ 的每个极小素理想都以此方式唯一地来自一个极小素理想
$\mathfrak p \subset A$（把 $A'$ 写成 $A$ 上的幂级数环即可容易证明
所有这些陈述）。因此，回顾极小素理想对应于不可约分支，下述定义有意义。

\begin{definition}
\label{stacks-geometry-definition-formal-branches}
设 $\mathcal{X}$ 是局部 Noether 概形 $S$ 上的局部有限型代数叠。
% P12-ERR-0153: see ../control/ERRATA_CANDIDATES.jsonl
设 $x_0 : \Spec(k) \to \mathcal{X}$ 是一态射，其中 $k$ 是 $S$ 上的
有限型域。{\it $\mathcal{X}$ 穿过 $x_0$ 的形式分支}，是任取
$\mathcal{X}$ 在 $x_0$ 处的一个半通用环时 $\Spec(A)$ 的不可约分支
所成的集合；对 $A$ 的不同选择，按上述步骤作认同。
% P12-ERR-0154: see ../control/ERRATA_CANDIDATES.jsonl
\end{definition}

\noindent
假设在定义 \ref{stacks-geometry-definition-formal-branches} 的情形中给定有限扩张
$l/k$。令 $x_{l, 0} : \Spec(l) \to \mathcal{X}$ 等于
$\Spec(l) \to \Spec(k)$ 与 $x_0$ 的复合。设
$A \to A'$ 如引理 \ref{stacks-geometry-lemma-versal-ring-field-extension} 中所述。
由于 $A \to A'$ 忠实平坦，态射
$$
\Spec(A') \to \Spec(A)
$$
把不可约分支（的一般点）映到不可约分支（的一般点）。
这将是一个满射，但一般并非双射。换言之，得到满射
$$
\text{形式分支：}\mathcal{X}\text{，穿过 }x_{l, 0}
\longrightarrow
\text{形式分支：}\mathcal{X}\text{，穿过 }x_0
$$
事实表明，若 $l/k$ 纯不可分，则该映射也单射
（如果将来需要，我们会在此补上精确陈述与证明）。

\begin{lemma}
\label{stacks-geometry-lemma-branches}
在定义 \ref{stacks-geometry-definition-formal-branches} 的情形中，存在一个典范满射：
从 $\mathcal{X}$ 穿过 $x_0$ 的形式分支集合，映到
$|\mathcal{X}|$ 中包含 $x_0$ 的 $|\mathcal{X}|$ 的不可约分支集合。
\end{lemma}

\begin{proof}
设 $A$ 如定义 \ref{stacks-geometry-definition-formal-branches} 中所述，并设
$\Spec(A) \to \mathcal{X}$ 如注 \ref{stacks-geometry-remark-upgrade} 中所述。
我们断言，$\Spec(A)$ 的一个不可约分支的一般点，映到
$|\mathcal{X}|$ 的一个不可约分支的一般点。取概形 $U$ 和满的光滑态射
$U \to \mathcal{X}$。考虑图表
$$
\xymatrix{
\Spec(A) \times_\mathcal{X} U \ar[d]_p \ar[r]_-q & U \ar[d]^f \\
\Spec(A) \ar[r]^j & \mathcal{X}
}
$$
由引理 \ref{stacks-geometry-lemma-versal-ring-flat}，$j$ 平坦，故 $q$ 平坦。
另一方面，$f$ 满且光滑，因而 $p$ 满且光滑。这说明不可约分支的任一
一般点 $\eta \in \Spec(A)$，都是代数空间
$\Spec(A) \times_\mathcal{X} U$ 的一个余维 $0$ 的点 $\eta'$ 的像
（记号见《空间的性质》第
\ref{spaces-properties-section-generic-points} 节，并在 \'etale
局部环上使用降链性）。由于 $q$ 平坦，$q(\eta')$ 是 $U$ 的余维 $0$
的点（论证相同）。由于 $U$ 是概形，$q(\eta')$ 是 $U$ 的一个
不可约分支的一般点。于是由引理 \ref{stacks-geometry-lemma-map-of-components}，
$q(\eta')$ 在 $|\mathcal{X}|$ 中的像的闭包是不可约分支，断言得证。

\medskip\noindent
显然，该断言给出了定义所需映射的机制。为证明其满射性，取
$u_0 \in U$，使其在 $|\mathcal{X}|$ 中映到 $x_0$。取 $u_0$ 的一个
仿射开邻域 $U' \subset U$。缩小 $U'$ 后，可以设 $U'$ 的每个
不可约分支都经过 $u_0$。于是可把 $\mathcal{X}$ 替换为
$|U'| \to |\mathcal{X}|$ 的像所对应的开子叠。因此可以设 $U$ 仿射，
% P12-ERR-0155: see ../control/ERRATA_CANDIDATES.jsonl
有一点 $u_0$ 映到 $x_0 \in |\mathcal{X}|$，并且 $U$ 的每个
不可约分支都经过 $u_0$。由《叠的性质》中的引理
\ref{stacks-properties-lemma-points-cartesian}，存在一点
$t \in |\Spec(A) \times_\mathcal{X} U|$，它映到 $\Spec(A)$ 的闭点以及
$u_0$。对平坦局部环同态使用降链性
$$
A \longrightarrow
\mathcal{O}_{\Spec(A) \times_\mathcal{X} U, \overline{t}}
\longleftarrow \mathcal{O}_{U, u_0}
$$
可知，$\mathcal{O}_{U, u_0}$ 的每个极小素理想，都是中间局部环的
一个极小素理想的像，而这样的极小素理想又映到 $A$ 的一个极小素理想。
这就证明了满射性。这里略去若干细节。
\end{proof}

\noindent
设 $A$ 是 Noether 完备局部环。于是 $\Spec(A)$ 的不可约分支具有重数；
见第 \ref{stacks-geometry-section-multiplicities} 节的引言。若
$A' = A[[t_1, \ldots, t_r]]$，则态射 $\Spec(A') \to \Spec(A)$
在不可约分支上诱导保持重数的双射（略去容易的证明）。
这一事实与定义 \ref{stacks-geometry-definition-formal-branches} 之前的讨论说明，
下述定义有意义。

\begin{definition}
\label{stacks-geometry-definition-multiplicity-formal-branches}
设 $\mathcal{X}$ 是局部 Noether 概形 $S$ 上的局部有限型代数叠。
% P12-ERR-0156: see ../control/ERRATA_CANDIDATES.jsonl
设 $x_0 : \Spec(k) \to \mathcal{X}$ 是一态射，其中 $k$ 是 $S$ 上的
有限型域。{\it $\mathcal{X}$ 穿过 $x_0$ 的一个形式分支的重数}，
是任取 $\mathcal{X}$ 在 $x_0$ 处的一个半通用环时，$\Spec(A)$ 的
相应不可约分支的重数（见上文讨论）。
\end{definition}

\begin{lemma}
\label{stacks-geometry-lemma-branches-multiplicity}
设 $\mathcal{X}$ 是局部 Noether 概形 $S$ 上的局部有限型代数叠。
% P12-ERR-0157: see ../control/ERRATA_CANDIDATES.jsonl
设 $x_0 : \Spec(k) \to \mathcal{X}$ 是一态射，其中 $k$ 是 $S$ 上的
有限型域，其像为 $s \in S$。若 $\mathcal{O}_{S, s}$ 是 G-环，
则引理 \ref{stacks-geometry-lemma-branches} 中的映射保持重数。
\end{lemma}

\begin{proof}
由引理 \ref{stacks-geometry-lemma-Artin-approximation-by-smooth-morphism}，可以设存在
一个光滑态射 $U \to \mathcal{X}$，其中 $U$ 是概形，并且 $U$ 有一个
$k$-值点 $u_0$，使得 $\mathcal{O}_{U, u_0}^\wedge$ 是
$\mathcal{X}$ 在 $x_0$ 处的半通用环。由引理
\ref{stacks-geometry-lemma-branches} 的证明中对映射的构造
（因为 $A = \mathcal{O}_{U, u_0}^\wedge$，此处大为简化），
只需证明：$U$ 的一个经过 $u_0$ 的不可约分支之重数，等于
$\Spec(\mathcal{O}_{U, u_0}^\wedge)$ 的映入它的任一不可约分支之重数。

\medskip\noindent
转化为交换代数，得到如下问题。设 $C = \mathcal{O}_{U, u_0}$。
它本质上是 $\mathcal{O}_{S, s}$ 上的有限型代数，因而是 G-环
（《代数的进一步讨论》中的命题
\ref{more-algebra-proposition-finite-type-over-G-ring}）。
于是 $A = C^\wedge$，所以 $C \to A$ 是正则环同态。
设 $\mathfrak q \subset C$ 是极小素理想，并设
$\mathfrak p \subset A$ 是位于 $\mathfrak q$ 上方的极小素理想。则
% P12-ERR-0158: see ../control/ERRATA_CANDIDATES.jsonl
$$
R = C_\mathfrak p \longrightarrow A_\mathfrak p = R'
$$
是 Artin 局部环之间的正则环同态。对这种环同态总有
$$
\text{长度}_R R = \text{长度}_{R'} R'
$$
这正是需要证明的，因为左边是 $U$ 上该分支的重数，右边是
$\Spec(A)$ 上该分支的重数。为看出这一等式，先由《代数》中的引理
\ref{algebra-lemma-pullback-module} 得
$$
\text{长度}_R(R) \text{长度}_{R'}(R'/\mathfrak m_R R') =
\text{长度}_{R'}(R')
$$
因此只需证明 $\mathfrak m_R R' = \mathfrak m_{R'}$；这是零维局部环
之间的正则同态的一个推论。
\end{proof}



\section{代数叠的维数理论}
\label{stacks-geometry-section-dimension-of-algebraic-stacks}

\noindent
就我们所知，文献中关于代数叠维数理论的主要结果出自
\cite{Osserman}，其中研究了余维和相对维数的概念。
% P12-ERR-0159: see ../control/ERRATA_CANDIDATES.jsonl
我们更细致地考察代数叠在一点处的维数，并证明若干结果，把一个态射在
源中某点处的纤维维数，与其源和目标的维数联系起来。我们还证明一项结果
（即下文引理 \ref{stacks-geometry-lemma-dimension-formula}），它使我们能够在适当假设下，
用半通用环计算代数叠在一点处的维数。
% P12-ERR-0160: see ../control/ERRATA_CANDIDATES.jsonl

\medskip\noindent
尽管我们并未总是试图把结果优化到最强形式，但大体上力求避免不必要的
假设。然而，在某些把代数叠的性质与它在一点处的半通用环之性质相比较的
结果中，我们只考虑如下代数叠：它在某个局部 Noether 概形基上局部有限
表示，并且该基的所有局部环都是 $G$-环。这样便可使用 Artin 逼近，
把半通用环的几何与叠本身的几何相比较。不过，我们证明的各项陈述未必都
真正需要这一限制性假设。由于我们所考虑的应用均满足它，在有帮助时采用
这一假设即可。

\medskip\noindent
若 $X$ 是概形，则把 $X$ 的维数 $\dim(X)$ 定义为 $X$ 的底拓扑空间的
Krull 维数；若 $x$ 是 $X$ 的一点，则把 $X$ 在 $x$ 处的维数
$\dim_x (X)$ 定义为所有包含 $x$ 的 $X$ 的开子集 $U$ 的维数之最小值；
见《性质》中的定义 \ref{properties-definition-dimension}。
有关系式 $\dim(X) = \sup_{x \in X} \dim_x(X)$；见《性质》中的引理
\ref{properties-lemma-dimension}。若 $X$ 局部 Noether，则
$\dim_x(X)$ 等于 $X$ 的所有经过 $x$ 的不可约分支在 $x$ 处维数的上确界。

\medskip\noindent
若 $X$ 是代数空间且 $x \in |X|$，则定义
$\dim_x X = \dim_u U,$，其中 $U$ 是任一具有 \'etale 满射
$U \to X$ 的概形，而 $u\in U$ 是位于 $x$ 上方的任一点；见
《空间的性质》中的定义
\ref{spaces-properties-definition-dimension-at-point}。令
$\dim(X) = \sup_{x \in |X|} \dim_x(X)$；见《空间的性质》中的定义
\ref{spaces-properties-definition-dimension}。

\begin{remark}
\label{stacks-geometry-remark-dimension-algebraic-space}
一般而言，代数空间 $X$ 在一点 $x$ 处的维数，未必等于其底拓扑空间
$|X|$ 在 $x$ 处的维数。例如，若 $k$ 是特征为零的域，且
$X =  \mathbf{A}^1_k / \mathbf{Z}$，则 $X$ 在其每一点处的维数都是
$1$（即 $\mathbf{A}^1_k$ 的维数）；而 $|X|$ 具有平凡拓扑，故其
Krull 维数为零。另一方面，《代数空间》中的例
\ref{spaces-example-infinite-product} 给出一个代数空间的例子：
它在每一点处的维数都是 $0$，但 $|X|$ 不可约、Krull 维数为 $1$，
并且具有一般点（所以 $|X|$ 在任一点处的维数都是 $1$）。
另见《空间的性质》第 \ref{spaces-properties-section-dimension} 节
对该例的讨论。

\medskip\noindent
另一方面，若 $X$ 是《合宜空间》中的定义
\ref{decent-spaces-definition-very-reasonable} 所称的
{\it 合宜的（decent）}代数空间（特别地，若 $X$ 拟分离；见
《合宜空间》第 \ref{decent-spaces-section-reasonable-decent} 节），
那么 $X$ 在 $x$ 处的维数实际上确实等于 $|X|$ 在 $x$ 处的维数；
见《合宜空间》中的引理
\ref{decent-spaces-lemma-dimension-decent-space}。
\end{remark}

\noindent
为了定义代数叠的维数，先引入如下相对维数概念会很有用：一个源为
代数空间、目标为代数叠的态射，在源中某点处的相对维数。这个定义稍显
繁复，仅仅是因为（与概形情形不同）代数叠或代数空间的点不能描述成
从一个域的谱出发的态射，而只能描述成这类态射的等价类。

\begin{definition}
\label{stacks-geometry-definition-relative-dimension}
若 $f : T \to \mathcal{X}$ 是从代数空间到代数叠的一个局部有限型态射，
% P12-ERR-0161: see ../control/ERRATA_CANDIDATES.jsonl
且 $t \in |T|$ 是一点，其像为 $x \in | \mathcal{X}|$，则如下定义
$f$ 在 $t$ 处的{\it 相对维数}，记作 $\dim_t(T_x),$：
选择一个代表 $x$ 的态射 $\Spec k \to \mathcal{X}$，其源为域的谱；
并选择一点 $t' \in |T \times_{\mathcal{X}} \Spec k|$，它在到
$|T|$ 的投影下映到 $t$（由《叠的性质》中的引理
\ref{stacks-properties-lemma-points-cartesian}，这样的点 $t'$ 存在）；
则
$$
\dim_t(T_x) = \dim_{t'}(T \times_{\mathcal{X}} \Spec k ).
$$
\end{definition}

\noindent
注意，由于 $T$ 是代数空间而 $\mathcal{X}$ 是代数叠，纤维积
$T \times_{\mathcal{X}} \Spec k$ 是代数空间，所以上述拟议定义右端的
量确实有定义（见上文讨论）。

\begin{remark}
\label{stacks-geometry-remark-relative-dimension}
(1)
容易验证 $\dim_t(T_x)$ 良定义，且与计算时所作选择无关。例如可使用
概形之间局部有限型态射的相对维数在基变换下不变这一事实；例如见
《态射》中的引理
\ref{morphisms-lemma-dimension-fibre-after-base-change}。

\medskip\noindent
(2)
当 $\mathcal{X}$ 也是代数空间时，直接检验即可确认这一定义与
《空间的态射》中的定义
\ref{spaces-morphisms-definition-dimension-fibre} 所给出的相对维数
定义一致。
\end{remark}

\noindent
接下来回顾下述引理；我们对局部 Noether 代数叠维数的研究以它为基础。

\begin{lemma}
\label{stacks-geometry-lemma-behaviour-of-dimensions-wrt-smooth-morphisms}
若 $f: U \to X$ 是局部 Noether 代数空间之间的光滑态射，且
$u \in |U|$ 的像为 $x \in |X|$，则
$$
\dim_u (U) = \dim_x(X) + \dim_{u} (U_x)
$$
其中 $\dim_u (U_x)$ 由定义 \ref{stacks-geometry-definition-relative-dimension} 给出。
\end{lemma}

\begin{proof}
见《空间的态射》中的引理
\ref{spaces-morphisms-lemma-smoothness-dimension-spaces}；注意，由注
\ref{stacks-geometry-remark-relative-dimension} (2)，此处使用的 $\dim_u (U_x)$
之定义与该处使用的定义一致。
\end{proof}

\begin{lemma}
\label{stacks-geometry-lemma-dimension-for-stacks}
% P12-ERR-0162: see ../control/ERRATA_CANDIDATES.jsonl
若 $\mathcal{X}$ 是局部 Noether 代数叠，且 $x \in |\mathcal{X}|$。
设 $U \to \mathcal{X}$ 是从代数空间到 $\mathcal{X}$ 的光滑态射，
并设 $u$ 是 $|U|$ 中映到 $x$ 的任一点。则
$$
\dim_x(\mathcal{X}) =  \dim_u(U) - \dim_{u}(U_x)
$$
其中相对维数 $\dim_u(U_x)$ 由定义
\ref{stacks-geometry-definition-relative-dimension} 给出，而 $\mathcal{X}$ 在 $x$
处的维数如《叠的性质》中的定义
\ref{stacks-properties-definition-dimension-at-point} 所述。
\end{lemma}

\begin{proof}
引理 \ref{stacks-geometry-lemma-behaviour-of-dimensions-wrt-smooth-morphisms} 可用来验证，
右端 $\dim_u(U) + \dim_u(U_x)$ 与光滑态射 $U \to \mathcal{X}$ 以及
$u \in |U|$ 的选择无关。
% P12-ERR-0163: see ../control/ERRATA_CANDIDATES.jsonl
略去细节。特别地，可以设 $U$ 是概形。在这种情形下，可以选取
$x$ 的代表为复合 $\Spec \kappa(u) \to U \to \mathcal{X}$ 来计算
$\dim_u(U_x)$，其中第一支态射是像为 $u \in U$ 的典范态射。
然后，若记 $R = U \times_{\mathcal{X}} U$，并用
$e : U \to R$ 表示对角态射，则相对维数在基变换下不变说明
$\dim_u(U_x) = \dim_{e(u)}(R_u)$。因此右端等于
$\dim_u (U) - \dim_{e(u)}(R_u) = \dim_x(\mathcal{X})$，正是所需。
\end{proof}

\begin{remark}
\label{stacks-geometry-remark-dimension-DM}
对于适当合宜的 Deligne--Mumford 叠（例如拟分离叠），
$\dim_x(\mathcal{X})$ 仍将等于拓扑定义的量
$\dim_x |\mathcal{X}|$。然而，对更一般的 Artin 叠，通常并非如此。
例如，若 $\mathcal{X} = [\mathbf{A}^1/\mathbf{G}_m]$
（在某个域上，并就 $\mathbf{G}_m$ 在 $\mathbf{A}^1$ 上的通常乘法
作用取商），则 $|\mathcal{X}|$ 有两个点，其中一点是另一点的特殊化
（对应于 $\mathbf{G}_m$ 在 $\mathbf{A}^1$ 上的两个轨道），故作为
拓扑空间其维数为 $1$；但对两个点 $x \in |\mathcal{X}|$，都有
$\dim_x (\mathcal{X}) = 0$。（一个更极端的例子是分类空间
$[\Spec k/\mathbf{G}_m]$，它在其唯一一点处的维数等于 $-1$。）
\end{remark}

\noindent
现在可以把定义 \ref{stacks-geometry-definition-relative-dimension} 推广到
（局部 Noether）代数叠之间的（局部有限型）态射。

\begin{definition}
\label{stacks-geometry-definition-relative-dimension-for-stacks}
若 $f : \mathcal{T} \to \mathcal{X}$ 是局部 Noether 代数叠之间的
局部有限型态射，且 $t \in |\mathcal{T}|$ 是一点，其像为
$x \in |\mathcal{X}|$，则如下定义 $f$ 在 $t$ 处的{\it 相对维数}，
记作 $\dim_t(\mathcal{T}_x),$：
选择一个代表 $x$ 的态射 $\Spec k \to \mathcal{X}$，其源为域的谱；
并选择一点 $t' \in |\mathcal{T} \times_{\mathcal{X}} \Spec k|$，
它在到 $|\mathcal{T}|$ 的投影下映到 $t$
（由《叠的性质》中的引理
\ref{stacks-properties-lemma-points-cartesian}，这样的点 $t'$ 存在；
% P12-ERR-0164: see ../control/ERRATA_CANDIDATES.jsonl
则
$$
\dim_t(\mathcal{T}_x) = \dim_{t'}(\mathcal{T} \times_{\mathcal{X}} \Spec k ).
$$
\end{definition}

\noindent
注意，由于 $\mathcal{T}$ 是代数叠而 $\mathcal{X}$ 是代数叠，纤维积
$\mathcal{T}\times_{\mathcal{X}} \Spec k$ 是代数叠，并由《叠的态射》
中的引理
\ref{stacks-morphisms-lemma-locally-finite-type-locally-noetherian}
知其局部 Noether。因此，上述拟议定义右端的量由《叠的性质》中的定义
\ref{stacks-properties-definition-dimension-at-point} 给出。

\begin{remark}
\label{stacks-geometry-remark-dimension-tangent-space-well-defined}
标准操作表明 $\dim_t(\mathcal{T}_x)$ 良定义，且与计算时所作选择无关。
\end{remark}

\noindent
现在建立相对维数的一些基本性质；它们是概形态射情形相应陈述的显然推广。

\begin{lemma}
\label{stacks-geometry-lemma-base-change-invariance-of-relative-dimension}
设有局部 Noether 叠的态射所成的笛卡儿方块
$$
\xymatrix{
\mathcal{T}' \ar[d]\ar[r] & \mathcal{T} \ar[d] \\
\mathcal{X}' \ar[r] & \mathcal{X}
}
$$
其中竖直态射局部有限型。若 $t' \in |\mathcal{T}'|$ 在
$|\mathcal{T}|$、$|\mathcal{X}'|$ 和 $|\mathcal{X}|$ 中的像分别为
$t$、$x'$ 和 $x$，则
$\dim_{t'}(\mathcal{T}'_{x'}) = \dim_{t}(\mathcal{T}_x).$
\end{lemma}

\begin{proof}
按定义，两边都可计算为同一个纤维积的维数。
\end{proof}

\begin{lemma}
\label{stacks-geometry-lemma-behaviour-of-dimensions-wrt-smooth-morphisms-stacky}
若 $f: \mathcal{U} \to \mathcal{X}$ 是局部 Noether 代数叠之间的
光滑态射，且 $u \in |\mathcal{U}|$ 的像为 $x \in |\mathcal{X}|$，则
$$
\dim_u (\mathcal{U}) = \dim_x(\mathcal{X}) + \dim_{u} (\mathcal{U}_x).
$$
\end{lemma}

\begin{proof}
取源为概形的光滑满射 $V \to \mathcal{U}$，并设 $v\in |V|$ 是映到
$u$ 的一点。复合 $V \to \mathcal{U} \to \mathcal{X}$ 也光滑，并由
引理 \ref{stacks-geometry-lemma-behaviour-of-dimensions-wrt-smooth-morphisms} 有
$\dim_x(\mathcal{X}) = \dim_v(V) - \dim_v(V_x),$
而 $\dim_u(\mathcal{U}) = \dim_v(V) - \dim_v(V_u).$ 因此
$$
\dim_u(\mathcal{U}) - \dim_x(\mathcal{X}) = \dim_v (V_x) - \dim_v (V_u).
$$

\medskip\noindent
选择 $x$ 的一个代表 $\Spec k \to \mathcal{X}$，并选择位于 $v$ 上方的
一点 $v' \in | V \times_{\mathcal{X}} \Spec k|$，其在
$|\mathcal{U}\times_{\mathcal{X}} \Spec k|$ 中的像为 $u'$；则按定义
$\dim_u(\mathcal{U}_x) = \dim_{u'}(\mathcal{U}\times_{\mathcal{X}} \Spec k),$
并且
$\dim_v(V_x) = \dim_{v'}(V\times_{\mathcal{X}} \Spec k).$

\medskip\noindent
现在
$V\times_{\mathcal{X}} \Spec k \to \mathcal{U}\times_{\mathcal{X}}\Spec k$
是光滑满射（因为它是这种态射的基变换），而其源是代数空间
（因为 $V$ 和 $\Spec k$ 是概形，$\mathcal{X}$ 是代数叠）。
因此再次按定义有
\begin{align*}
\dim_{u'}(\mathcal{U}\times_{\mathcal{X}} \Spec k)
& =
\dim_{v'}(V\times_{\mathcal{X}} \Spec k) -
\dim_{v'}(V \times_{\mathcal{X}} \Spec k)_{u'}) \\
& = \dim_v(V_x) -
\dim_{v'}( (V\times_{\mathcal{X}} \Spec k)_{u'}).
\end{align*}
又
$V\times_{\mathcal{X}} \Spec k \cong
V\times_{\mathcal{U}} (\mathcal{U}\times_{\mathcal{X}} \Spec k),$
所以引理 \ref{stacks-geometry-lemma-base-change-invariance-of-relative-dimension} 表明
$\dim_{v'}((V\times_{\mathcal{X}} \Spec k)_{u'})  = \dim_v(V_u).$
综合以上各式，得到
$$
\dim_u(\mathcal{U}) - \dim_x(\mathcal{X}) =
\dim_u(\mathcal{U}_x),
$$
正是所需。
\end{proof}

\begin{lemma}
\label{stacks-geometry-lemma-relative-dimension-is-semi-continuous}
设 $f: \mathcal{T} \to \mathcal{X}$ 是代数叠之间的局部有限型态射。
% P12-ERR-0165: see ../control/ERRATA_CANDIDATES.jsonl
\begin{enumerate}
\item
函数 $t \mapsto \dim_t(\mathcal{T}_{f(t)})$ 在 $|\mathcal{T}|$ 上
上半连续。
\item 若 $f$ 光滑，则函数
$t \mapsto \dim_t(\mathcal{T}_{f(t)})$ 在 $|\mathcal{T}|$ 上局部常值。
\end{enumerate}
\end{lemma}

\begin{proof}
先假设 $\mathcal{T}$ 是概形 $T$。设 $U \to \mathcal{X}$ 是源为概形的
光滑满射，并令 $T' = T \times_{\mathcal{X}} U$。设
$f': T' \to U$ 是 $f$ 在 $U$ 上的拉回，并设 $g: T' \to T$ 为投影。

\medskip\noindent
引理 \ref{stacks-geometry-lemma-base-change-invariance-of-relative-dimension} 表明，
对 $t' \in T'$ 有
$\dim_{t'}(T'_{f'(t')}) = \dim_{g(t')}(T_{f(g(t'))}),$。
由于 $g$ 光滑且满（它是光滑满射的基变换），$g$ 在底拓扑空间上
诱导的映射由《空间的性质》中的引理
\ref{spaces-properties-lemma-topology-points} 连续且开，并且满射。
因此，只需指出：态射 $f'$ 的第 (1) 项由《空间的态射》中的引理
\ref{spaces-morphisms-lemma-openness-bounded-dimension-fibres} 得出，
而第 (2) 项可由《态射》中的引理
\ref{morphisms-lemma-flat-finite-presentation-CM-fibres-relative-dimension}
或《态射》中的引理
\ref{morphisms-lemma-smooth-omega-finite-locally-free} 得出
（二者都给出概形情形的结果，由此可完全仿照《空间的态射》中的引理
\ref{spaces-morphisms-lemma-openness-bounded-dimension-fibres} 推出
代数空间的相应结果。
% P12-ERR-0166: see ../control/ERRATA_CANDIDATES.jsonl

\medskip\noindent
现在回到一般情形，并选择源为概形的光滑满射
$h:V \to \mathcal{T}$。若 $v \in V$，则基本上按定义有
$$
\dim_{h(v)}(\mathcal{T}_{f(h(v))}) =
\dim_{v}(V_{f(h(v))}) - \dim_{v}(V_{h(v)}).
$$
由于 $V$ 是概形，我们已经证明该等式右边第一项上半连续
（若 $f$ 光滑，它甚至局部常值），而第二项事实上局部常值。
所以它们的差上半连续（若 $f$ 光滑，则局部常值），因而函数
$\dim_{h(v)}(\mathcal{T}_{f(h(v))})$ 在 $|V|$ 上上半连续
（若 $f$ 光滑，则局部常值）。由于态射
$|V| \to |\mathcal{T}|$ 开且满，引理得证。
\end{proof}

\noindent
在继续展开理论之前，先证明两个与概形维数理论有关的引理。

\medskip\noindent
为说明第一个引理的背景，注意，若 $X$ 是有限维概形，则由于
$\dim X$ 被定义为各维数 $\dim_x X$ 的上确界，存在一点
$x \in X$，使得 $\dim_x X = \dim X$。下一个引理表明，还可以取
$x$ 为有限型点。

\begin{lemma}
\label{stacks-geometry-lemma-dimension-achieved-by-finite-type-point}
若 $X$ 是有限维概形，则存在闭点（从而是有限型点）$x \in X$，使得
$\dim_x X = \dim X$。
\end{lemma}

\begin{proof}
令 $d = \dim X$，并取 $X$ 的不可约闭子集的一条极大严格递降链，记为
\begin{equation}
\label{stacks-geometry-equation-maximal-chain}
Z_0 \supset Z_1 \supset \ldots \supset Z_d.
\end{equation}
子集 $Z_d$ 是 $X$ 的极小不可约闭子集，故 $Z_d$ 的每一点都是
$Z_d$ 的一般点。由于概形 $X$ 的底拓扑空间是清醒的，可知 $Z_d$
是单点集，仅由一个闭点 $x \in X$ 构成。若 $U$ 是 $x$ 的任一邻域，
则链
$$
U\cap Z_0 \supset U\cap Z_1 \supset \ldots \supset U\cap Z_d = Z_d =
\{x\}
$$
是 $U$ 的不可约闭子集的一条严格递降链，说明 $\dim U \geq d$。
因此 $\dim_x X \geq d$。反向不等式显然，引理得证。
\end{proof}

\noindent
下一个引理说明，在满足若干温和附加假设的不可约概形上，
$\dim_x X$ 是{\it 常值}函数。

\begin{lemma}
\label{stacks-geometry-lemma-constancy-of-dimension}
若 $X$ 是有限维的不可约、Jacobson、链状且局部 Noether 概形，
则对 $X$ 的每个非空开子集 $U$ 都有 $\dim U = \dim X$。
等价地，$\dim_x X$ 是 $X$ 上的常值函数。
\end{lemma}

\begin{proof}
两个断言的等价性直接由定义得出。于是，设 $U\subset X$ 是非空开子集。
显然 $\dim U \leq \dim X$，而需要证明 $\dim U \geq \dim X.$
令 $d = \dim X$，并取 $X$ 的不可约闭子集的一条极大严格递降链，记为
$$
X = Z_0 \supset Z_1 \supset \ldots \supset Z_d.
$$
由于 $X$ 是 Jacobson 的，极小不可约闭子集 $Z_d$ 等于某个闭点 $x$
所成的 $\{x\}$。

\medskip\noindent
若 $x \in U,$，则
$$
U = U \cap Z_0  \supset U\cap Z_1 \supset \ldots \supset
U\cap Z_d = \{x\}
$$
是 $U$ 的不可约闭子集的一条严格递降链，故
$\dim U \geq d$，正是所需。因此可设 $x \not\in U.$

\medskip\noindent
考虑平坦态射 $\Spec \mathcal{O}_{X,x} \to X$。$X$ 的非空
（从而稠密）开子集 $U$ 拉回为开子集
$V \subset \Spec \mathcal{O}_{X,x}$。把 $U$ 替换为一个非空拟紧、
从而 Noether 的开子集后，可以设嵌入 $U \to X$ 是拟紧态射。
由于拟紧态射的概形论像之形成与平坦基变换交换，
% P12-ERR-0167: see ../control/ERRATA_CANDIDATES.jsonl
由《态射》中的引理
\ref{morphisms-lemma-flat-base-change-scheme-theoretic-image} 可知
$V$ 在 $\Spec \mathcal{O}_{X,x}$ 中稠密，因而特别非空；当然
$x \not\in V.$（这里也用 $x$ 表示 $\Spec \mathcal{O}_{X,x}$ 的闭点，
因为其像等于给定点 $x \in X$。）现在
$\Spec \mathcal{O}_{X,x} \setminus \{x\}$ 是 Jacobson 的；
% P12-ERR-0168: see ../control/ERRATA_CANDIDATES.jsonl
见《性质》中的引理
\ref{properties-lemma-complement-closed-point-Jacobson}。所以 $V$ 包含
$\Spec \mathcal{O}_{X,x} \setminus \{x\}$ 的一个闭点 $z$。$z$ 的像
在 $X$ 中的闭包，是 $X$ 的一个包含 $x$ 的不可约闭子集 $Z$；
它与 $U$ 的交非空，并且不存在真包含于 $Z$ 而又真包含
$\{x\}$ 的不可约闭子集（因为拉回到 $\Spec \mathcal{O}_{X,x}$
在 $X$ 中包含 $x$ 的不可约闭子集，与
$\Spec \mathcal{O}_{X,x}$ 的不可约闭子集之间诱导双射）。
由于 $U \cap Z$ 是 $U$ 的非空闭子集，它包含一个在 $X$ 中闭的点
$u$（因为 $X$ 是 Jacobson 的）。又因为 $U\cap Z$ 是不可约集合 $Z$
的非空（从而稠密）开子集，而它包含一个不属于 $U$ 的点，即 $x$，
所以嵌入 $\{u\} \subset U\cap Z$ 是真包含。

\medskip\noindent
由于 $X$ 链状，链
$$
X = Z_0 \supset Z \supset \{x\} = Z_d
$$
可细化为长度 $d+1$ 的链，而该链必具有形式
$$
X = Z_0 \supset W_1 \supset \ldots \supset W_{d-1} = Z \supset \{x\} = Z_d.
$$
由于 $U\cap Z$ 非空，遂得到
$$
U = U \cap Z_0 \supset U \cap W_1\supset \ldots \supset U\cap W_{d-1}
= U\cap Z \supset \{u\}
$$
这是 $U$ 的不可约闭子集的一条长度为 $d+1$ 的严格递降链，说明
$\dim U \geq d$，正是所需。
\end{proof}

\noindent
引理 \ref{stacks-geometry-lemma-constancy-of-dimension} 的叠论类比将在下文引理
\ref{stacks-geometry-lemma-irreducible-implies-equidimensional} 中证明；但在此之前，
必须引入一个附加定义。之所以需要它，是因为概形为链状这一概念并非
\'etale 局部的（见《代数》中的注
\ref{algebra-remark-universally-catenary-does-not-descend} 的例子；
这使得很难定义代数空间或代数叠为链状是什么意思
（见 \cite[page 3]{Osserman} 中的讨论）。
% P12-ERR-0169: see ../control/ERRATA_CANDIDATES.jsonl
对维数理论的某些方面，下述定义似乎很好地替代了所缺少的链状代数叠概念。

\begin{definition}
\label{stacks-geometry-definition-pseudo-catenary}
称局部 Noether 代数叠 $\mathcal{X}$ 是{\it 拟链状的}，如果存在一个
光滑满射 $U \to \mathcal{X}$，其源是万有链状概形。
\end{definition}

\begin{example}
\label{stacks-geometry-example-pseudo-catenary}
若 $\mathcal{X}$ 在万有链状局部 Noether 概形 $S$ 上局部有限型，且
$U\to \mathcal{X}$ 是源为概形的光滑满射，则复合
$U \to \mathcal{X} \to S$ 局部有限型，所以 $U$ 万有链状；
% P12-ERR-0170: see ../control/ERRATA_CANDIDATES.jsonl
见《态射》中的引理
\ref{morphisms-lemma-universally-catenary-local}。因此
$\mathcal{X}$ 拟链状。
\end{example}

\noindent
下述引理表明，拟链状性质沿有限型态射传递。

\begin{lemma}
\label{stacks-geometry-lemma-catenary-covers}
若 $\mathcal{X}$ 是拟链状局部 Noether 代数叠，而
$\mathcal{Y} \to \mathcal{X}$ 是局部有限型态射，
% P12-ERR-0171: see ../control/ERRATA_CANDIDATES.jsonl
则存在一个光滑满射 $V \to \mathcal{Y}$，其源是万有链状概形；因此
$\mathcal{Y}$ 也拟链状。
\end{lemma}

\begin{proof}
根据假设，可以找到源为万有链状概形的光滑满射
$U \to \mathcal{X}$。基变换
$U\times_{\mathcal{X}} \mathcal{Y}$ 是代数叠；取源为概形的光滑满射
$V \to U\times_{\mathcal{X}} \mathcal{Y}$。于是复合
$V \to U\times_{\mathcal{X}} \mathcal{Y} \to \mathcal{Y}$ 光滑且满
（因为它是光滑满射的复合），而态射
$V \to U\times_{\mathcal{X}} \mathcal{Y} \to U$ 局部有限型
（因为它是局部有限型态射的复合）。由于 $U$ 万有链状，由《态射》中的
引理 \ref{morphisms-lemma-universally-catenary-local}，$V$ 万有链状，
正如所断言。
\end{proof}

\noindent
现在研究函数 $\dim_x(\mathcal{X})$ 在 $|\mathcal{X}|$ 上关于
$|\mathcal{X}|$ 的不可约分支的行为，以及若干相关问题；这里
$\mathcal{X}$ 是某个局部 Noether 叠。

\begin{lemma}
\label{stacks-geometry-lemma-irreducible-implies-equidimensional}
若 $\mathcal{X}$ 是 Jacobson、拟链状且局部 Noether 的代数叠，
并且 $|\mathcal{X}|$ 不可约，则 $\dim_x(\mathcal{X})$ 是
$|\mathcal{X}|$ 上的常值函数。
\end{lemma}

\begin{proof}
只需证明 $\dim_x(\mathcal{X})$ 在 $|\mathcal{X}|$ 上局部常值；
这样它必然常值（因为 $|\mathcal{X}|$ 不可约，故连通）。
由于 $\mathcal{X}$ 拟链状，可以找到光滑满射
$U \to \mathcal{X}$，其中 $U$ 是万有链状概形。若 $\{U_i\}$ 是
$U$ 的拟紧开子概形覆盖，则可把 $U$ 替换为 $\coprod U_i,$；
只需证明函数 $u \mapsto \dim_{f(u)}(\mathcal{X})$ 在 $U_i$ 上
局部常值。由于每次只对一个 $U_i$ 检验，现在略去下标，直接写作
$U$ 而非 $U_i$。由于 $U$ 拟紧，它是有限多个不可约分支的并，
记为 $T_1 \cup \ldots \cup T_n$。注意，每个 $T_i$ 都是 Jacobson、
链状且局部 Noether 的，因为它是 Jacobson、链状且局部 Noether
概形 $U$ 的闭子概形。

\medskip\noindent
由引理 \ref{stacks-geometry-lemma-behaviour-of-dimensions-wrt-smooth-morphisms} 有
$\dim_{f(u)}(\mathcal{X}) = \dim_{u}(U) - \dim_{u}(U_{f(u)}).$
引理 \ref{stacks-geometry-lemma-relative-dimension-is-semi-continuous} (2) 表明，
由于 $f$ 光滑，右端第二项在 $U$ 上局部常值；因此必须证明
$\dim_u(U)$ 在 $U$ 上局部常值。由于 $\dim_u(U)$ 是所有包含 $u$ 的
$U$ 的分支 $T_i$ 的维数 $\dim_u T_i$ 中的最大值，只需证明：
若一点 $u$ 位于两个不同分支 $T_i$ 与 $T_j$ 上（其中 $i \neq j$），
则 $\dim_u T_i = \dim_u T_j$；再注意，在不可约、Jacobson、链状且
局部 Noether 的概形 $T$ 上，$t\mapsto \dim_t T$ 是常值函数
（这由引理 \ref{stacks-geometry-lemma-constancy-of-dimension} 得出）。

\medskip\noindent
令 $V = T_i \setminus (\bigcup_{i' \neq i} T_{i'})$，并令
$W = T_j \setminus (\bigcup_{i' \neq j} T_{i'})$。于是 $V$ 和 $W$
各自都是 $U$ 的非空开子集，故各自在 $|\mathcal{X}|$ 中具有非空开像。
由于 $|\mathcal{X}|$ 不可约，这两个非空开子集在
$|\mathcal{X}|$ 中有非空交。取位于该交中的一点 $x$，并取映到 $x$
的点 $v \in V$ 和 $w\in W$。于是有
$$
\dim T_i = \dim V = \dim_v (U) = \dim_x (\mathcal{X}) + \dim_v (U_x)
$$
并且类似地有
$$
\dim T_j = \dim W = \dim_w (U) = \dim_x (\mathcal{X}) + \dim_w (U_x).
$$
由于 $u \mapsto \dim_u (U_{f(u)})$ 在 $U$ 上局部常值，而
$T_i \cup T_j$ 连通（它是两个不可约、从而连通，并具有非空交的集合之并），
可知 $\dim_v (U_x) = \dim_w(U_x)$。比较前两个等式，即得
$\dim T_i = \dim T_j$，正是所需。
\end{proof}

\begin{lemma}
\label{stacks-geometry-lemma-closed-immersions}
若 $\mathcal{Z} \hookrightarrow \mathcal{X}$ 是局部 Noether 代数叠的
闭浸入，而 $z \in |\mathcal{Z}|$ 的像为
$x \in |\mathcal{X}|$，则
$\dim_z (\mathcal{Z}) \leq \dim_x(\mathcal{X})$。
\end{lemma}

\begin{proof}
取源为概形的光滑满射 $U\to \mathcal{X}$；基变换后的态射
$V = U\times_{\mathcal{X}} \mathcal{Z} \to \mathcal{Z}$ 也光滑且满，
而投影 $V \to U$ 是闭浸入。若 $v \in |V|$ 映到
$z \in |\mathcal{Z}|$，并以 $u$ 表示 $v$ 在 $|U|$ 中的像，则显然
$\dim_v(V) \leq \dim_u(U)$；另一方面，由引理
\ref{stacks-geometry-lemma-base-change-invariance-of-relative-dimension} 有
$\dim_v (V_z) = \dim_u(U_x)$。因此
$$
\dim_z(\mathcal{Z})  = \dim_v(V) - \dim_v(V_z)
\leq \dim_u(U) - \dim_u(U_x) = \dim_x(\mathcal{X}),
$$
正如所断言。
\end{proof}

\begin{lemma}
\label{stacks-geometry-lemma-dimension-via-components}
若 $\mathcal{X}$ 是局部 Noether 代数叠，且
$x \in |\mathcal{X}|$，则
$\dim_x(\mathcal{X}) = \sup_{\mathcal{T}} \{ \dim_x(\mathcal{T}) \} $，
其中 $\mathcal{T}$ 遍历 $|\mathcal{X}|$ 的所有经过 $x$ 的不可约分支
（赋予其诱导约化结构）。
\end{lemma}

\begin{proof}
引理 \ref{stacks-geometry-lemma-closed-immersions} 表明，对每个经过点 $x$ 的不可约分支
$\mathcal{T}$，都有
$\dim_x (\mathcal{T}) \leq \dim_x(\mathcal{X})$。因此为证明引理，
只需证明
\begin{equation}
\label{stacks-geometry-equation-desired-inequality}
\dim_x(\mathcal{X}) \leq
\sup_{\mathcal{T}} \{\dim_x(\mathcal{T})\}.
\end{equation}
设 $U\to\mathcal{X}$ 是由概形给出的光滑覆盖。若 $T$ 是 $U$ 的一个
不可约分支，则以 $\mathcal{T}$ 表示其像在 $\mathcal{X}$ 中的闭包；
它是 $\mathcal{X}$ 的不可约分支。设 $u \in U$ 是映到 $x$ 的一点。
于是
$\dim_x(\mathcal{X})=\dim_uU-\dim_uU_x=\sup_T\dim_uT-\dim_uU_x$，
其中上确界取遍 $U$ 的所有经过 $u$ 的不可约分支。选择一个达到该
上确界的分支 $T$，并注意
$\dim_x(\mathcal{T})=\dim_uT-\dim_u T_x$。
所需不等式（\ref{stacks-geometry-equation-desired-inequality}）现在由显然的不等式
$\dim_u T_x \leq \dim_u U_x.$ 得出。
（注意，若 $\Spec k \to \mathcal{X}$ 是 $x$ 的一个代表，则
$T\times_{\mathcal{X}} \Spec k$ 是
$U\times_{\mathcal{X}} \Spec k$ 的闭子空间。）
\end{proof}

\begin{lemma}
\label{stacks-geometry-lemma-dimension-at-finite-type-point}
若 $\mathcal{X}$ 是局部 Noether 代数叠，且
$x \in |\mathcal{X}|$，则对 $\mathcal{X}$ 的任一包含 $x$ 的开子叠
$\mathcal{V}$，存在有限型点 $x_0 \in |\mathcal{V}|$，使得
$\dim_{x_0}(\mathcal{X}) = \dim_x(\mathcal{V})$。
\end{lemma}

\begin{proof}
取源为概形的光滑满射 $f:U \to \mathcal{X}$，并考虑函数
$u \mapsto \dim_{f(u)}(\mathcal{X});$。由于 $f$ 诱导的态射
$|U| \to |\mathcal{X}|$ 是开的（因为 $f$ 光滑）且满的
（根据假设），并且把有限型点映到有限型点（这正是
$|\mathcal{X}|$ 的有限型点之定义），只需证明：对任意 $u \in U$
以及 $u$ 的任一开邻域，都存在该邻域中的有限型点 $u_0$，使得
$\dim_{f(u_0)}(\mathcal{X}) = \dim_{f(u)}(\mathcal{X}).$
在对问题作此重述后，不再需要 $f$ 满，故可把 $U$ 替换为所考虑点
$u$ 的开邻域，从而把问题约化为证明：对每个 $u \in U$，存在有限型点
$u_0 \in U$，使得
$\dim_{f(u_0)}(\mathcal{X}) = \dim_{f(u)}(\mathcal{X}).$
由引理 \ref{stacks-geometry-lemma-behaviour-of-dimensions-wrt-smooth-morphisms}，
$\dim_{f(u)}(\mathcal{X}) = \dim_u(U) - \dim_u(U_{f(u)}),$
而
$\dim_{f(u_0)}(\mathcal{X}) = \dim_{u_0}(U) - \dim_{u_0}(U_{f(u_0)}).$
由于 $f$ 光滑，当 $u_0$ 在 $U$ 上变化时，表达式
$\dim_{u_0}(U_{f(u_0)})$ 局部常值
（由引理 \ref{stacks-geometry-lemma-relative-dimension-is-semi-continuous} (2)）。
必要时进一步在 $u$ 附近缩小 $U$，可以设它常值。于是问题变成证明：
可以找到有限型点 $u_0 \in U$，使得
$\dim_{u_0}(U) = \dim_u(U)$。按定义，$\dim_u U$ 是各维数
$\dim V$ 的最小值，其中 $V$ 取遍这样的开邻域 $V$：它们位于
原概形中并包含 $u$；所以可进一步
在 $u$ 附近缩小 $U$，使得 $\dim_u U = \dim U$。所需点 $u_0$
的存在性遂由引理 \ref{stacks-geometry-lemma-dimension-achieved-by-finite-type-point}
得出。
\end{proof}

\begin{lemma}
\label{stacks-geometry-lemma-monomorphing-a-component-in-of-the-right-dimension}
设 $\mathcal{T} \hookrightarrow \mathcal{X}$ 是代数叠的局部有限型
单态射，其中 $\mathcal{X}$（从而 $\mathcal{T}$ 也）是 Jacobson、
拟链状且局部 Noether 的。
% P12-ERR-0172: see ../control/ERRATA_CANDIDATES.jsonl
再假设 $\mathcal{T}$ 不可约，具有某个（有限）维数 $d$，而
$\mathcal{X}$ 约化且维数小于或等于 $d$。则存在 $\mathcal{T}$ 的
非空开子叠 $\mathcal{V}$，使得诱导的单态射
$\mathcal{V} \hookrightarrow \mathcal{X}$ 是开浸入，并把
$\mathcal{V}$ 认同为 $\mathcal{X}$ 的一个不可约分支的开子集。
\end{lemma}

\begin{proof}
取源为概形的光滑满射 $f:U \to \mathcal{X}$；由于 $\mathcal{X}$
约化，$U$ 必然也约化。记
$U' = \mathcal{T}\times_{\mathcal{X}} U$。基变换后的态射
$U' \to U$ 是代数空间的局部有限型单态射，因此可表；
% P12-ERR-0173: see ../control/ERRATA_CANDIDATES.jsonl
见《空间的态射》中的引理
\ref{spaces-morphisms-lemma-locally-quasi-finite-separated-representable}
和
\ref{spaces-morphisms-lemma-monomorphism-loc-finite-type-loc-quasi-finite}。
由于 $U$ 是概形，$U'$ 也是概形。投影
$f': U' \to \mathcal{T}$ 仍是光滑满射。设 $u' \in U'$，其像为
$u \in U$。引理
\ref{stacks-geometry-lemma-base-change-invariance-of-relative-dimension} 表明
$\dim_{u'}(U'_{f(u')}) = \dim_u(U_{f(u)}),$；另一方面，由引理
\ref{stacks-geometry-lemma-irreducible-implies-equidimensional} 以及对
$\mathcal{T}$ 和 $\mathcal{X}$ 的假设，有
$\dim_{f'(u')}(\mathcal{T}) =d
\geq \dim_{f(u)}(\mathcal{X})$。因此
\begin{equation}
\label{stacks-geometry-equation-dim-inequality}
\dim_{u'} (U') = \dim_{u'} (U'_{f(u')}) + \dim_{f'(u')}(\mathcal{T})
\\
\geq \dim_u (U_{f(u)}) + \dim_{f(u)}(\mathcal{X}) = \dim_u (U).
\end{equation}
由于 $U' \to U$ 是局部有限型单态射，特别地它非分歧；所以由非分歧
态射的 \'etale 局部结构，即《\'Etale 态射》中的引理
\ref{etale-lemma-finite-unramified-etale-local}，可找到交换图表
$$
\xymatrix{
V' \ar[r]\ar[d] & V \ar[d] \\
U' \ar[r] & U
}
$$
其中概形 $V'$ 非空，竖直箭头为 \'etale 态射，而上方水平箭头是闭浸入。
把 $V$ 替换为一个拟紧开子集，使其像与 $U'$ 的像有非空交，并把
$V'$ 替换为 $V$ 的原像，还可设 $V$（从而 $V'$）拟紧。
由于 $V$ 也局部 Noether，它遂为 Noether，因而是有限多个不可约分支的并。

\medskip\noindent
由于 \'etale 态射保持逐点维数，
% P12-ERR-0174: see ../control/ERRATA_CANDIDATES.jsonl
由《下降》中的引理 \ref{descent-lemma-dimension-at-point-local}
和（\ref{stacks-geometry-equation-dim-inequality}）可推出：对任一点
$v' \in V'$，若其像为 $v \in V$，则
$\dim_{v'}( V') \geq \dim_v(V)$。特别地，$V'$ 的像不可能包含于
$V$ 的两个不同不可约分支的交中；因此可以找到 $V$ 的至少一个
不可约开子集，它与 $V'$ 有非空交。把 $V$ 替换为此子集后，可以设
$V$ 是整概形（因为它既约化又不可约）。由上述维数不等式可知，
闭浸入 $V' \hookrightarrow V$ 实际上是同构。
若以 $W$ 表示 $V'$ 在 $U'$ 中的像，则 $W$ 是 $U'$ 的非空开子集
（因为 \'etale 态射是开映射），而诱导的单态射 $W \to U$ 为
\'etale 态射（因为它在源上 \'etale 局部，即拉回到 $V'$ 后如此），
故为开浸入（因为它是 \'etale 单态射）。因此，若以 $\mathcal{V}$
表示 $W$ 在 $\mathcal{T}$ 中的像，则 $\mathcal{V}$ 是
$\mathcal{T}$ 的稠密（等价地，非空）开子叠，其像在
$\mathcal{X}$ 的一个不可约分支中稠密。最后，注意态射
% P12-ERR-0175: see ../control/ERRATA_CANDIDATES.jsonl
$\mathcal{V} \to \mathcal{X}$ 光滑（因为它与光滑态射
$W\to \mathcal{V}$ 的复合光滑），并且也是单态射，故为开浸入。
\end{proof}

\begin{lemma}
\label{stacks-geometry-lemma-dims-of-images}
设 $f: \mathcal{T} \to \mathcal{X}$ 是 Jacobson、拟链状且局部
Noether 的代数叠之间的局部有限型态射，其源不可约，目标拟分离；
% P12-ERR-0176: see ../control/ERRATA_CANDIDATES.jsonl
并设 $\mathcal{Z} \hookrightarrow \mathcal{X}$ 表示
$\mathcal{T}$ 的概形论像。则对所有 $t \in |T|$，都有
% P12-ERR-0177: see ../control/ERRATA_CANDIDATES.jsonl
$\dim_t( \mathcal{T}_{f(t)}) \geq \dim \mathcal{T}  - \dim \mathcal{Z}$；
而且 $|\mathcal{T}|$ 有一个非空（等价地，稠密）开子集，使得其上等号成立。
\end{lemma}

\begin{proof}
用 $\mathcal{Z}$ 替换 $\mathcal{X}$ 后，可以并且确实设 $f$
在概形论意义下支配，并且 $\mathcal{X}$ 不可约。由纤维维数的
上半连续性（引理 \ref{stacks-geometry-lemma-relative-dimension-is-semi-continuous} (1)），
只需证明等式
$\dim_t( \mathcal{T}_{f(t)}) =\dim \mathcal{T}  - \dim \mathcal{Z}$
对位于 $\mathcal{T}$ 的某个非空开子叠中的 $t$ 成立。因此在论证中，
始终可以自由地把 $\mathcal{T}$ 替换为非空开子叠。

\medskip\noindent
设 $T' \to \mathcal{T}$ 是源为概形的光滑满射，并设 $T$ 是 $T'$ 的
非空拟紧开子集。由于 $\mathcal{Y}$ 拟分离，可知
% P12-ERR-0178: see ../control/ERRATA_CANDIDATES.jsonl
$T \to  \mathcal{Y}$ 拟紧（把《叠的态射》中的引理
\ref{stacks-morphisms-lemma-quasi-compact-permanence} 应用于态射
$T \to \mathcal{Y} \to \Spec \mathbf{Z}$）。因此，若把
$\mathcal{T}$ 替换为 $T$ 在 $\mathcal{T}$ 中的像，则可以设
（援引《叠的态射》中的引理
\ref{stacks-morphisms-lemma-surjection-from-quasi-compact}
% P12-ERR-0179: see ../control/ERRATA_CANDIDATES.jsonl
态射 $f:\mathcal{T} \to \mathcal{X}$ 拟紧。

\medskip\noindent
若取光滑满射 $U \to \mathcal{X}$，其中 $U$ 是概形，则引理
\ref{stacks-geometry-lemma-map-of-components} 保证可以找到 $U$ 的一个不可约开子集
$V$，使得 $V \to \mathcal{X}$ 光滑且在概形论意义下支配。
由于拟紧态射的概形论支配性在平坦基变换下保持，概形论支配态射 $f$
的基变换 $\mathcal{T} \times_{\mathcal{X}} V \to V$ 仍在概形论
意义下支配。设 $Z$ 是一个具有到该纤维积的光滑满射的概形；
于是 $Z \to \mathcal{T} \times_{\mathcal{X}} V \to V$
仍在概形论意义下支配。因此可找到 $Z$ 的一个不可约分支 $C$，
它在概形论意义下支配 $V$。由于复合
$Z \to \mathcal{T}\times_{\mathcal{X}} V \to \mathcal{T}$ 光滑，
而 $\mathcal{T}$ 不可约，引理 \ref{stacks-geometry-lemma-map-of-components} 表明，
源的任一不可约分支在 $|\mathcal{T}|$ 中具有稠密像。现在把 $C$
替换为一个与 $Z$ 的所有其他不可约分支不交的非空开子集 $W$，
再把 $\mathcal{T}$ 和 $\mathcal{X}$ 替换为 $W$ 和 $V$ 的像
（并应用引理 \ref{stacks-geometry-lemma-irreducible-implies-equidimensional}，
可知这不改变 $\mathcal{T}$ 或 $\mathcal{X}$ 的维数）。
若以 $\mathcal{W}$ 表示态射
$W \to \mathcal{T}\times_{\mathcal{X}} V$ 的像，则 $\mathcal{W}$
在 $\mathcal{T}\times_{\mathcal{X}} V$ 中开
（因为态射 $W \to \mathcal{T}\times_{\mathcal{X}} V$ 光滑），
并且不可约（因为它是不可约概形的像）。于是最终得到交换图表
$$
\xymatrix{
W \ar[dr] \ar[r]  & \mathcal{W} \ar[r] \ar[d] & V \ar[d] \\
& \mathcal{T} \ar[r] & \mathcal{X}
}
$$
其中 $W$ 和 $V$ 是概形，竖直箭头光滑且满，斜箭头以及左上水平箭头光滑，
% P12-ERR-0180: see ../control/ERRATA_CANDIDATES.jsonl
而诱导态射
$\mathcal{W} \to \mathcal{T}\times_{\mathcal{X}} V$ 是开浸入。
利用这一图表以及引理陈述中所涉各维数的定义，可把引理的验证约化到
已知的概形情形。

\medskip\noindent
固定 $w \in |W|$，并将其像依次记作
$w' \in |\mathcal{W}|$、$t \in |\mathcal{T}|$、
$v$（在 $|V|$ 中）和 $x$（在 $|\mathcal{X}|$ 中）。基本上按定义
（利用 $\mathcal{W}$ 在 $\mathcal{T}\times_{\mathcal{X}} V$ 中开，
以及基变换的纤维是纤维的基变换）得到等式
$$
\dim_v V_x = \dim_{w'} \mathcal{W}_t
$$
和
$$
\dim_t \mathcal{T}_x = \dim_{w'} \mathcal{W}_v.
$$
由引理 \ref{stacks-geometry-lemma-behaviour-of-dimensions-wrt-smooth-morphisms}
（图表中的斜箭头和右侧竖直箭头，把 $W$ 和 $V$ 实现为叠
$\mathcal{T}$ 和 $\mathcal{X}$ 的概形光滑覆盖），得到
$$
\dim_t \mathcal{T} = \dim_w W - \dim_w W_t
$$
以及
$$
\dim_x \mathcal{X} = \dim_v V - \dim_v V_x.
$$
合并这些等式，得到
$$
\dim_t \mathcal{T}_x - \dim_t \mathcal{T} + \dim_x \mathcal{X}
= \dim_{w'} \mathcal{W}_v - \dim_w W + \dim_w W_t + \dim_v V -
\dim_{w'} \mathcal{W}_t
$$
由于 $W \to \mathcal{W}$ 是光滑满射，在态射
$\Spec \kappa(v) \to V$ 上作基变换后仍然如此
（把 $W \to \mathcal{W}$ 视为 $V$ 上的态射）。
从这一光滑态射得到下列两个等式中的第一个：
$$
\dim_w W_v - \dim_{w'} \mathcal{W}_v = \dim_w (W_v)_{w'} = \dim_w W_{w'};
$$
第二个等式由直接比较所涉及的两个纤维得出。
类似地，若把 $W \to \mathcal{W}$ 视为 $\mathcal{T}$ 上的概形态射，
并在点 $t \in |\mathcal{T}|$ 的某个代表上作基变换，就得到等式
$$
\dim_w W_t - \dim_{w'} \mathcal{W}_t = \dim_w (W_t)_{w'} = \dim_w W_{w'}.
$$
综合以上各式，得到
$$
\dim_t \mathcal{T}_x - \dim_t \mathcal{T} + \dim_x \mathcal{X}
=  \dim_w W_v - \dim_w W + \dim_v V.
$$
我们的目标是证明，当 $t$ 取遍某个非空开子集时，该等式左边为零。
% P12-ERR-0181: see ../control/ERRATA_CANDIDATES.jsonl
当 $w$ 在 $W$ 的一个非空开子集上变化时，其像
$t \in |\mathcal{T}|$ 在 $|\mathcal{T}|$ 的一个非空开子集上变化
（因为 $W \to \mathcal{T}$ 光滑）。

\medskip\noindent
因此问题约化为证明：若 $W\to V$ 是不可约局部 Noether 概形之间的
局部有限型概形论支配态射，则存在其点的一个非空开子集，使得对
其中的 $w\in W$ 有
$\dim_w W_v =\dim_w W - \dim_v V$（其中 $v$ 表示 $w$ 在 $V$ 中的像）。
这是标准事实；为方便读者，回顾其证明。

\medskip\noindent
可把 $W$ 和 $V$ 替换为其底约化子概形，而不改变这一等式成立与否，
故可设它们事实上是整概形。由于 $\dim_w W_v$ 在 $W,$ 上局部常值，
% P12-ERR-0182: see ../control/ERRATA_CANDIDATES.jsonl
必要时把 $W$ 替换为非空开子集，可以设 $\dim_w W_v$ 常值，记其值为
$d$。选择该开子集为仿射的，还可设态射 $W\to V$ 实际上有限型。
必要时把 $V$ 替换为非空开子集
（随后在该开子集上拉回 $W$；所得拉回非空，因为拟紧概形论支配态射的
平坦基变换仍然概形论支配），还可设 $W$ 在 $V$ 上平坦。
于是态射 $W\to V$ 在《态射》中的定义
\ref{morphisms-definition-relative-dimension-d} 的意义下具有相对维数
$d$，并由《态射》中的引理
\ref{morphisms-lemma-rel-dimension-dimension} 得到
$\dim_w(W) = \dim_v(V) + d,$，正是所需。
\end{proof}

\begin{remark}
\label{stacks-geometry-remark-negative-dimension}
我们指出，在上一引理的情形中，未必有
$\dim \mathcal{T} \geq \dim \mathcal{Z}$；这并不与该引理陈述中的
不等式矛盾，因为态射 $f$ 的纤维仍是代数叠，因而维数可能为负。
例如，取 $k$ 为一个域，并把该引理应用于态射
$[\Spec k/\mathbf{G}_m] \to \Spec k$，即可看出这一点。

\medskip\noindent
若进一步假设该引理陈述中的态射 $f$ 是拟 DM 的
（其含义见《叠的态射》中的定义
\ref{stacks-morphisms-definition-separated}；例如，可由代数空间表示的
态射是拟 DM 的），则该态射在目标各点上的纤维都是拟 DM 代数叠，
因而具有非负维数。在这种情形下，该引理确实蕴含
$\dim \mathcal{T} \geq \dim \mathcal{Z}$。事实上，我们得到下面这个
更一般的结果。
\end{remark}

\begin{lemma}
\label{stacks-geometry-lemma-dims-of-images-two}
% P12-ERR-0183: see ../control/ERRATA_CANDIDATES.jsonl
设 $f: \mathcal{T} \to \mathcal{X}$ 是 Jacobson、拟链状且局部
Noether 的代数叠之间的局部有限型态射。假设它拟 DM，其源不可约，
其目标拟分离，并以
$\mathcal{Z} \hookrightarrow \mathcal{X}$ 表示 $\mathcal{T}$
的概形论像。则
$\dim \mathcal{Z} \leq \dim \mathcal{T}$，而且下列两个条件中
恰有一个成立：
\begin{enumerate}
\item 对每个有限型点 $t \in |T|,$ 都有
% P12-ERR-0184: see ../control/ERRATA_CANDIDATES.jsonl
$\dim_t(\mathcal{T}_{f(t)}) > 0,$ 在这种情形下，
$\dim \mathcal{Z} < \dim \mathcal{T}$；或者
\item $\mathcal{T}$ 和 $\mathcal{Z}$ 具有相同维数。
\end{enumerate}
\end{lemma}

\begin{proof}
正如上一备注所指出的，拟 DM 叠的维数总是非负的。因此，对所有
$t \in |\mathcal{T}|$ 都有
$\dim_t \mathcal{T}_{f(t)} \geq 0$，并且等式
$$
\dim_t \mathcal{T}_{f(t)} = \dim_t \mathcal{T} - \dim_{f(t)} \mathcal{Z}
$$
对一个稠密开子集中的点 $t\in |\mathcal{T}|$ 成立。
\end{proof}





\section{局部环的维数}
\label{stacks-geometry-section-dimension-local-ring}

\noindent
代数叠在通常意义下其实并没有局部环，但可以如下定义局部环的维数。

\begin{lemma}
\label{stacks-geometry-lemma-dimension-local-ring-pre}
设 $\mathcal{X}$ 是局部 Noether 代数叠。设 $U \to \mathcal{X}$
是光滑态射，并设 $u \in U$。则
$$
\dim(\mathcal{O}_{U, \overline{u}}) -
\dim(\mathcal{O}_{R_u, e(\overline{u})}) =
2\dim(\mathcal{O}_{U, \overline{u}}) -
\dim(\mathcal{O}_{R, e(\overline{u})})
$$
这里 $R = U \times_\mathcal{X} U$，投影为 $s, t : R \to U$，
对角态射为 $e : U \to R$，而 $R_u$ 是 $s : R \to U$ 在 $u$
上的纤维。
\end{lemma}

\begin{proof}
这是因为
$s : \mathcal{O}_{U, \overline{u}} \to \mathcal{O}_{R, e(\overline{u})}$
是 Noether 局部环之间的平坦局部同态，因而
$$
\dim(\mathcal{O}_{R, e(\overline{u})}) =
\dim(\mathcal{O}_{U, \overline{u}}) +
\dim(\mathcal{O}_{R_u, e(\overline{u})})
$$
由《代数》中的引理
\ref{algebra-lemma-dimension-base-fibre-equals-total} 成立。
\end{proof}

\begin{lemma}
\label{stacks-geometry-lemma-dimension-local-ring}
设 $\mathcal{X}$ 是局部 Noether 代数叠。
% P12-ERR-0185: see ../control/ERRATA_CANDIDATES.jsonl
设 $x \in |\mathcal{X}|$ 是有限型点
（见《叠的态射》中的定义
\ref{stacks-morphisms-definition-finite-type-point}）。
设 $d \in \mathbf{Z}$。下列条件等价：
\begin{enumerate}
\item 存在一个概形 $U$、一个光滑态射 $U \to \mathcal{X}$，以及
映到 $x$ 的有限型点 $u \in U$，使得
$2\dim(\mathcal{O}_{U, \overline{u}}) -
\dim(\mathcal{O}_{R, e(\overline{u})}) = d$；并且
\item 对任意概形 $U$、任意光滑态射 $U \to \mathcal{X}$，以及
任意映到 $x$ 的有限型点 $u \in U$，都有
$2\dim(\mathcal{O}_{U, \overline{u}}) -
\dim(\mathcal{O}_{R, e(\overline{u})}) = d$。
\end{enumerate}
这里 $R = U \times_\mathcal{X} U$，投影为 $s, t : R \to U$，
对角态射为 $e : U \to R$，而 $R_u$ 是 $s : R \to U$ 在 $u$
上的纤维。
\end{lemma}

\begin{proof}
假设有 $x$ 的两个光滑邻域 $(U, u)$ 和 $(U', u')$，其中 $u$ 和
$u'$ 都是有限型点。缩小 $U$ 和 $U'$ 后，可设 $u$ 和 $u'$ 都是
闭点（这由有限型点的定义得出）。然后选取满 \'etale 态射
$W \to U \times_\mathcal{X} U'$。设 $W_u$ 为 $W \to U$ 在 $u$
上的纤维，并设 $W_{u'}$ 为 $W \to U'$ 在 $u'$ 上的纤维。
由于 $u$ 和 $u'$ 映到 $|\mathcal{X}|$ 的同一点，
$W_u \cap W_{u'}$ 非空。因此可选取一个闭点 $w \in W$，它同时映到
$u$ 和 $u'$。这把问题约化为下一段中的讨论。

\medskip\noindent
设 $(U', u') \to (U, u)$ 是 $x$ 的两个光滑邻域之间的光滑态射，
且 $u$ 和 $u'$ 都是闭点。目标是证明：用 $(U, u)$ 定义的不变量
与用 $(U', u')$ 定义的不变量相同。为此注意，
$\mathcal{O}_{U, u} \to \mathcal{O}_{U', u'}$
是 Noether 局部环之间的平坦局部同态，因而
$$
\dim(\mathcal{O}_{U', \overline{u}'}) =
\dim(\mathcal{O}_{U, \overline{u}}) +
\dim(\mathcal{O}_{U'_u, \overline{u}'})
$$
由《代数》中的引理
\ref{algebra-lemma-dimension-base-fibre-equals-total} 成立。
（这里略去把局部环的性质同其严格 Hensel 化的性质联系起来所需的
全部步骤；见《代数进阶》中的章节
\ref{more-algebra-section-permanence-henselization}。）
另一方面，有
$$
R' = U' \times_{U, t} R \times_{s, U} U'
$$
因此可见
$$
\dim(\mathcal{O}_{R', e(\overline{u}')}) =
\dim(\mathcal{O}_{R, e(\overline{u})}) +
\dim(\mathcal{O}_{U'_u \times_u U'_u, (\overline{u}', \overline{u}')})
$$
要证明该引理，只需证明
$$
\dim(\mathcal{O}_{U'_u \times_u U'_u, (\overline{u}', \overline{u}')}) =
2\dim(\mathcal{O}_{U'_u, \overline{u}'})
$$
注意这并不总成立（例如，若 $U'_u$ 是一条曲线，而 $u'$ 是这条曲线的
一般点）。不过，我们知道 $u'$ 是在 $u$ 上局部有限型的代数空间
$U'_u$ 的闭点。在这种情形下，等式成立：首先，由
《簇》中的引理
\ref{varieties-lemma-dimension-product-locally-algebraic}，
有
$\dim_{(u', u')}(U'_u \times_u U'_u) = 2\dim_{u'}(U'_u)$；
其次，在局部代数概形的闭点处，维数与局部环的维数一致，见《簇》中的
引理 \ref{varieties-lemma-dimension-locally-algebraic}。
我们略去把这些概形结论转述为代数空间语言的过程。
\end{proof}

\begin{definition}
\label{stacks-geometry-definition-dimension-local-ring}
设 $\mathcal{X}$ 是局部 Noether 代数叠，并设
$x \in |\mathcal{X}|$ 是有限型点。若引理
\ref{stacks-geometry-lemma-dimension-local-ring} 中的等价条件成立，则称其中的
$d \in \mathbf{Z}$ 为
{\it $\mathcal{X}$ 在 $x$ 处的局部环的维数}。
\end{definition}

\noindent
准确地说，这一定义由引理 \ref{stacks-geometry-lemma-dimension-local-ring-pre}
和《叠的性质》中的定义
\ref{stacks-properties-definition-dimension-at-point} 所启发。
% P12-ERR-0186: see ../control/ERRATA_CANDIDATES.jsonl
本节最后建立一个公式，以便用 $\mathcal{X}$ 在 $x$ 处的半通用环的
性质计算 $\dim_x(\mathcal{X})$。

\begin{lemma}
\label{stacks-geometry-lemma-dimension-formula}
设 $\mathcal{X}$ 是局部 Noether 概形 $S$ 上的局部有限型代数叠。
设 $x_0 : \Spec(k) \to \mathcal{X}$ 是一个态射，其中 $k$ 是
$S$ 上的有限型域。按备注 \ref{stacks-geometry-remark-groupoid-defo}，用一个
余群胚 $(A, B, s, t, c)$ 表示
$\mathcal{F}_{\mathcal{X}, k, x_0}$；这里的两个环都是剩余域为
$k$ 的 Noether 完备局部 $S$-代数。则
$$
\text{局部环维数：}\mathcal{X}\text{ 在 }x_0 =
2\dim A - \dim B
$$
\end{lemma}

\begin{proof}
设 $s \in S$ 是 $x_0$ 的像。若 $\mathcal{O}_{S, s}$ 是 G-环
（这是实践中几乎总能满足的条件），则可如下证明该引理。
由引理 \ref{stacks-geometry-lemma-Artin-approximation-by-smooth-morphism}，可以找到
一个光滑态射 $U \to \mathcal{X}$，其源是概形，并且包含一个剩余域为
$k$ 的点 $u_0 \in U$，使得诱导态射
% P12-ERR-0187: see ../control/ERRATA_CANDIDATES.jsonl
$\Spec(k) \to U \to \mathcal{X}$ 与 $x_0$ 重合，且
$A = \mathcal{O}_{U, u_0}^\wedge$。记
$R = U \times_\mathcal{X} U$。于是可把
$\mathcal{O}_{R, e(u_0)}^\wedge$ 与 $B$ 认同。所需等式遂由定义得出。

\medskip\noindent
在证明余下部分中，我们说明如何在一般情形下证明该引理，但建议读者
跳过这部分。

\medskip\noindent
先证明右边与 $(A, B, s, t, c)$ 的选择无关。具体地，设
$(A', B', s', t', c')$ 是第二种选择。由于 $A$ 和 $A'$ 都是
$\mathcal{X}$ 在 $x_0$ 处的半通用环，必要时交换 $A$ 和 $A'$ 后，
可以选取一个与 $A$ 和 $A'$ 上给定的半通用形式对象
$\xi$ 和 $\xi'$ 相容的形式光滑映射 $A \to A'$。
回忆 $\widehat{\mathcal{C}}_\Lambda$ 有余积，而且这些余积由
$\Lambda$ 上的完备张量积给出；见《形式形变理论》中的引理
\ref{formal-defos-lemma-CLambdahat-coproducts}。
那么，$B$ pro-表示这样一个函子：它把 $\xi$ 推前到
$A \widehat{\otimes}_\Lambda A$ 后所得的两个对象之间的同构。
$B'$ 的情形类似。由此得到
$$
B' =
B \otimes_{(A \widehat{\otimes}_\Lambda A)}
(A' \widehat{\otimes}_\Lambda A')
$$
不难看出
$$
A \widehat{\otimes}_\Lambda A \longrightarrow
A \widehat{\otimes}_\Lambda A' \longrightarrow
A' \widehat{\otimes}_\Lambda A'
$$
是形式光滑的，其相对维数等于形式光滑映射 $A \to A'$ 的相对维数的
$2$ 倍。（这可由一般原理得出；也可直接看出，因为在此特定情形中，
$A'$ 是 $A$ 上关于 $r$ 个变量的幂级数环。）因此 $B \to B'$
形式光滑，其相对维数为 $2(\dim(A') - \dim(A))$，正是所需。

\medskip\noindent
接着，设 $l/k$ 是有限扩张。
% P12-ERR-0188: see ../control/ERRATA_CANDIDATES.jsonl
设 $x_{l, 0} : \Spec(l) \to \mathcal{X}$ 为诱导的点。我们断言，
该公式右边对 $x_0$ 和对 $x_{l, 0}$ 取相同的值。
这可以通过按引理 \ref{stacks-geometry-lemma-versal-ring-field-extension} 选取
$A \to A'$，并完全照搬上一段的论证来证明。细节从略。

\medskip\noindent
最后，照引理 \ref{stacks-geometry-lemma-versal-ring-flat} 的证明进行论证，可以利用
前两段中的相容性，把问题约化为第一段所讨论的情形：此时 $A$ 是某个
在 $\mathcal{X}$ 上光滑的概形 $U$ 在 $u_0$ 处的完备局部环，而
% P12-ERR-0189: see ../control/ERRATA_CANDIDATES.jsonl
$u_0$ 是有限型点。细节从略。
\end{proof}
